\textbf{(D) Proof of Theorems 2 and 3.}

The Lie algebra homomorphism \(\beta \circ \alpha: \mathrm{L}_{\mathbf{Z}}(\mathrm{X}) \to \mathrm{A}(\mathrm{X})\) restricted to \(\mathbf{X}\) is the identity by (3) and hence is the canonical injection (§ 3, no. 1). Therefore \(\alpha\) is injective and hence bijective by (B); this proves Theorem 3. As \(\beta \circ \alpha = f \circ \varepsilon \circ \alpha\) is injective and \(\alpha\) is bijective, \(\varepsilon\) is injective. For all \(n \geqslant 1\) ,

\[
\varepsilon_ {n} \colon \mathrm{C} ^ {n} / \mathrm{C} ^ {n + 1} \rightarrow \mathrm{F} ^ {n} / \mathrm{F} ^ {n + 1}
\]

is injective and hence

\[
\mathrm{C} ^ {n} \cap \mathrm{F} ^ {n + 1} = \mathrm{C} ^ {n + 1}.
\]

\(C^{1} = F = F^{1}\) ; if \(C^{n} = F^{n}\) , then \(C^{n} \cap F^{n+1} = F^{n+1}\) whence \(C^{n+1} = F^{n+1}\) which proves Theorem 2 by induction on \(n \geqslant 1\) .

COROLLARY.

\[
\bigcap_ {n \geqslant 1} \mathrm{C} ^ {n} \mathrm{F} (\mathrm{X}) = \{e \}.
\]

Applying Theorem 2 with K = Z and r = 0,

\[
\bigcap_ {n \geqslant 1} \mathrm{C} ^ {n} \mathrm{F} (\mathrm{X}) = \bigcap_ {n \geqslant 1} \bar {g} (1 + \hat {\mathrm{A}} _ {n} (\mathrm{X})) = \bar {g} \left(\bigcap_ {n \geqslant 1} (1 + \hat {\mathrm{A}} _ {n} (\mathrm{X}))\right) = \bar {g} (1) = \{e \}.
\]

Remark. Let H be a Hall set relative to X (§2, no. 10). Let M be the magma defined by the law of composition \((x,y)\mapsto (x,y) = x^{-1}y^{-1}xy\) on F(X) and let \(\phi\) be the homomorphism of M(X) into M whose restriction to X is the identity. The elements of \(\phi (\mathrm{H})\) are called the basic commutators of F(X) associated with the Hall set H. For every integer \(n\geqslant 1\) , let \(\mathrm{H}_{\mathfrak{n}}\) be the subset of H consisting of the elements of length \(n\) ; we know (§2, no. 11, Theorem 1) that the canonical mapping of \(\mathrm{H}_{\mathfrak{n}}\) into \(\mathbf{L}_{\mathbf{Z}}(\mathbf{X})\) is a basis of the Abelian group \(\mathbf{L}_{\mathbf{Z}}^{n} (\mathbf{X})\) . Moreover, \(\phi (\mathrm{H_n})\subset \mathrm{C^n}\) ; for all \(m\in \mathrm{H}_{\mathfrak{n}}\) , let \(\phi_{\mathfrak{n}}(m)\) denote the class mod. \(\mathrm{C}^{n + 1}\) of \(\phi (m)\in \mathrm{C}^n\) . Theorem 3 then shows that \(\phi_{\mathfrak{n}}\) is a bijection of \(\mathrm{H}_{\mathfrak{n}}\) onto a basis of the Abelian group \(\mathrm{C}^n /\mathrm{C}^{n + 1}\) . It follows immediately that, for all \(w\in F(\mathbf{X})\) and all \(i\geqslant 1\) , there exists a unique element \(\alpha_{i}\) of \(\mathbf{Z}^{(\mathrm{H}_{i})}\) such that, for \(n\geqslant 1\) ,

\[
w = \prod_ {i = 1} ^ {n} \prod_ {m \in H _ {i}} \phi (m) ^ {\alpha_ {i} (m)} \mod C ^ {n + 1},\tag{4}
\]

where the product is calculated according to the total ordering given on H.

Example. Suppose that X is a set with two elements x, y and let \(H_{1} = \{x, y\}\) , \(H_{2} = \{xy\}\) . Every element w of F(X) can therefore be written

\[
w \equiv x ^ {a} y ^ {b} (x, y) ^ {c} \bmod . \mathrm{C} ^ {3} \quad \text { with } a, b, c \text { in } \mathbf {Z}.
\]

For \(w = (xy)^n\) , \(a = b = n\) and \(c = n(1 - n)/2\) (cf.~Exercise 9), whence

\[
(x y) ^ {n} \equiv x ^ {n} y ^ {n} (x, y) ^ {n (1 - n) / 2} \mod C ^ {3}.
\]

\subsection{p-FILTRATION OF FREE GROUPS}

In this no., \(p\) denotes a prime number and we assume that \(\mathbf{K} = \mathbf{F}_p\) . Let \(\gamma\) be the homomorphism of \(\mathrm{F}(\mathbf{X})\) into \(\Gamma(\mathbf{X})\) defined by \(\gamma(x) = 1 + x\) for \(x\) in \(\mathbf{X}\) ; we write \(\mathrm{F}_n^{(p)}(\mathbf{X}) = \gamma(1 + \hat{\Lambda}_n(\mathbf{X}))\) . The sequence \((\mathrm{F}_n^{(p)}(\mathbf{X}))_{n\geq 1}\) is an integral central filtration on \(\mathrm{F}(\mathbf{X})\) , which is separated since \(\gamma\) is injective (no. 3, Theorem 1). It is called the \(p\) -filtration on \(\mathrm{F}(\mathbf{X})\) .

PROPOSITION 2. Suppose that X is finite. For every integer \(n \geqslant 1\) , the group \(F(X) / F_n^{(p)}(X)\) is a finite \(p\) -group of nilpotency class \(\leqslant n\) .

Arguing by induction on \(n\) , it suffices to prove that \(\mathrm{F}_{n}^{(p)}(\mathbf{X}) / \mathrm{F}_{n + 1}^{(p)}(\mathbf{X})\) is a finite commutative \(p\) -group for all \(n \geqslant 1\) . For all \(w \in \mathrm{F}_{n}^{(p)}(\mathbf{X})\) , the element \(\gamma(w) - 1\) of \(\hat{\mathbf{A}}(\mathbf{X})\) is of order \(\geqslant n\) ; we denote by \(\delta_n(w)\) the homogeneous component of \(\gamma(w) - 1\) of degree \(n\) . The mapping \(\delta_n: \mathrm{F}_{n}^{(p)}(\mathbf{X}) \to \mathrm{A}^n (\mathbf{X})\) is a homomorphism with kernel \(\mathrm{F}_{n + 1}^{(p)}(\mathbf{X})\) (§ 4, no. 5, Proposition 3) and hence \(\mathrm{F}_{n}^{(p)}(\mathbf{X}) / \mathrm{F}_{n + 1}^{(p)}(\mathbf{X})\) is isomorphic to a subgroup of \(\mathrm{A}^n (\mathbf{X})\) . Since \(\mathbf{X}\) is finite, \(\mathrm{A}^n (\mathbf{X})\) is a finite-dimensional vector space over \(\mathbf{F}_p\) and hence a finite commutative \(p\) -group and so is \(\mathrm{F}_{n}^{(p)}(\mathbf{X}) / \mathrm{F}_{n + 1}^{(p)}(\mathbf{X})\) .

PROPOSITION 3. For all \(w \neq 1\) in \(\mathrm{F}(\mathbf{X})\) , there exist a finite \(p\) -group \(\mathbf{G}\) and a homomorphism \(f\) of \(\mathrm{F}(\mathbf{X})\) into \(\mathbf{G}\) such that \(f(w) \neq 1\) .

There exist elements \(x_{1},\ldots ,x_{r}\) of \(\mathbf{X}\) and integers \(n_1,\dots ,n_r\) such that \(w = x_1^{n_1}\dots x_r^{n_r}\) . Let \(\mathrm{Y} = \{x_1,\dots,x_r\}\) . The canonical injection of \(\mathrm{Y}\) into \(\mathbf{X}\) extends to a homomorphism \(\alpha : \mathrm{F}(\mathrm{Y})\to \mathrm{F}(\mathrm{X})\) ; on the other hand, let \(\beta\) be the homomorphism of \(\mathrm{F}(\mathrm{X})\) into \(\mathrm{F}(\mathrm{Y})\) whose restriction to \(\mathbf{Y}\) is the identity and which maps \(\mathbf{X} - \mathbf{Y}\) to \(\{1\}\) . Then \(\beta (\alpha (y)) = y\) for \(y\in \mathbf{Y}\) and hence \(\beta \circ \alpha\) is the identity automorphism of \(\mathrm{F}(\mathrm{Y})\) . There obviously exists \(w^{\prime}\) in \(\mathrm{F}(\mathrm{Y})\) such that \(w = \alpha (w^{\prime})\) ; then \(\beta (w) = w^{\prime}\neq 1\) ; now \(\bigcap_{n\geqslant 1}\mathbf{F}_n^{(p)}(\mathbf{Y}) = \{1\}\) and there therefore exists an integer \(n\geqslant 1\) such that \(\beta (w)\notin \mathbf{F}_{\mathrm{n}}^{(p)}(\mathbf{Y})\) . By Proposition 2, the group \(\mathbf{G} = \mathbf{F}(\mathbf{Y}) / \mathbf{F}_{\mathrm{n}}^{(p)}(\mathbf{Y})\) is a finite \(p\) -group. If \(f\) is the composition of \(\beta\) and the canonical homomorphism of \(\mathrm{F}(\mathrm{Y})\) onto \(\mathbf{G}\) , then \(f(w)\neq 1\) .

COROLLARY. The intersection of the normal subgroups of finite index in F(X) is \{1\}.

\section{THE HAUSDORFF SERIES}

In this paragraph we assume that K is a field of characteristic 0.

\subsection{EXPONENTIAL AND LOGARITHM IN FILTERED ALGEBRAS}

Let \(\mathbf{A}\) be a unital associative algebra which is Hausdorff and complete under a real filtration \((\mathbf{A}_{\alpha})\) . We write \(m = A_0^+ = \bigcup_{\alpha > 0} A_\alpha\) .

For \(x \in \mathfrak{m}\) , the family \((x^n / n!)_{n \in \mathbf{N}}\) is summable. We write

\[
e ^ {x} = \exp x = \sum_ {n \geqslant 0} x ^ {n} / n!.\tag{1}
\]

Then \(\exp (x)\in 1 + \mathfrak{m}\) and the mapping \(\exp \colon \mathfrak{m}\to 1 + \mathfrak{m}\) is called the exponential mapping of A.

For all \(y \in 1 + \mathfrak{m}\) , the family \(((-1)^{n-1}(y - 1)^n / n)_{n \geqslant 1}\) is summable. We write

\[
\log y = \sum_ {n \geqslant 1} (- 1) ^ {n - 1} (y - 1) ^ {n} / n.\tag{2}
\]

Then \(\log y \in m\) and the mapping \(\log: 1 + m \to m\) is called the logarithmic mapping of A.

PROPOSITION 1. The exponential mapping is a homeomorphism of \(\mathfrak{m}\) onto \(1 + \mathfrak{m}\) and the logarithmic mapping is the inverse homeomorphism.

For \(x \in \mathrm{A}_{\alpha}, \frac{x^n}{n!} \in \mathrm{A}_{n\alpha}\) . It follows that the series defining the exponential converges uniformly on each of the sets \(\mathrm{A}_{\alpha}\) for \(\alpha > 0\) ; as \(\mathrm{A}_{\alpha}\) is open in \(m\) and \(m = \bigcup_{\alpha > 0} \mathrm{A}_{\alpha}\) , the exponential mapping is continuous. It can be similarly shown that the logarithmic mapping is continuous.

Let \(e\) and \(l\) be the formal power series with no constant term

\[
e (\mathrm{X}) = \sum_ {n \geqslant 1} \frac {\mathrm{X} ^ {n}}{n !}, \quad l (\mathrm{X}) = \sum_ {n \geqslant 1} (- 1) ^ {n - 1} \mathrm{X} ^ {n} / n.
\]

We know (Algebra, Chapter IV, § 6, no. 9) that \(e(l(\mathbf{X})) = l(e(\mathbf{X})) = \mathbf{X}\) on \(\hat{\mathrm{A}}(\{\mathbf{X}\}) = \mathrm{K}[[\mathbf{X}]]\) . By substitution (§ 5, no. 1), we deduce that

\[
e (l (x)) = l (e (x)) = x
\]

for \(x\in \mathfrak{m}\) ; as

\[
\exp x = e (x) + 1, \quad \log (1 + x) = l (x)
\]

it follows immediately that

\[
\log \exp x = x, \quad \exp \log (1 + x) = 1 + x
\]

for \(x\) in \(\mathfrak{m}\) , whence the proposition.

Remarks. (1) If \(x \in \mathfrak{m}\) , \(y \in \mathfrak{m}\) and \(x\) and \(y\) commute, then

\[
\exp (x + y) = \exp (x) \exp (y),
\]

since the family \(\left(\frac{x^i}{i!}\cdot \frac{y^j}{j!}\right)_{i,j\in \mathbf{N}}\) is summable (cf.~Algebra, Chapter IV, § 6, no. 9, Proposition 11).

\begin{enumerate}
\def\labelenumi{(\arabic{enumi})}
\setcounter{enumi}{1}
\item
  As the series \(e\) and \(l\) are without constant term and \(A_{\alpha}\) is a closed ideal of \(A\) , \(\exp A_{\alpha} \subset 1 + A_{\alpha}\) and \(\log (1 + A_{\alpha}) \subset A_{\alpha}\) whence \(\exp A_{\alpha} = 1 + A_{\alpha}\) and \(\log (1 + A_{\alpha}) = A_{\alpha}\) for \(\alpha > 0\) .
\item
  Let B be a complete Hausdorff filtered unital associative algebra and \(n = \bigcup_{\alpha > 0} B_{\alpha}\) . Let f be a continuous unital homomorphism of A into B such that \(f(\mathfrak{m}) \subset \mathfrak{n}\) . Then \(f(\exp x) = \exp f(x)\) for \(x \in \mathfrak{m}\) and \(f(\log y) = \log f(y)\) for \(y \in 1 + \mathfrak{m}\) ; we show for example the first of these formulae:
\end{enumerate}

\[
f (\exp x) = \sum_ {n \geqslant 0} f (x ^ {n}) / n! = \sum_ {n \geqslant 0} f (x) ^ {n} / n! = \exp f (x).
\]

\begin{enumerate}
\def\labelenumi{(\arabic{enumi})}
\setcounter{enumi}{3}
\tightlist
\item
  Let E be a unital associative algebra. If \(a\) is a nilpotent element of E, the family \(\left(\frac{a^n}{n!}\right)_{n\in \mathbf{N}}\) has finite support and we write \(\exp a = \sum_{n\geqslant 0}a^n /n!\) . An element \(b\) is unipotent if \(b - 1\) is nilpotent; then we write
\end{enumerate}

\[
\log b = \sum_ {n \geqslant 1} (- 1) ^ {n - 1} (b - 1) ^ {n} / n.
\]

We deduce from the relations \(e(l(\mathbf{X})) = l(e(\mathbf{X})) = \mathbf{X}\) that the mapping \(a \mapsto \exp a\) is a bijection of the set of nilpotent elements of E onto the set of unipotent elements of E and that \(b \mapsto \log b\) is the inverse mapping.

\subsection{HAUSDORFF GROUP}

Let \(\mathbf{X}\) be a set. We use the notation of § 5, nos. 1 and 2. The free Lie algebra \(\mathrm{L}(\mathbf{X})\) is identified with its canonical image in \(\mathrm{A}(\mathbf{X})\) (§ 3, no. 1, Theorem 1). We denote by \(\hat{\mathrm{L}}(\mathbf{X})\) the closure of \(\mathrm{L}(\mathbf{X})\) in \(\hat{\mathrm{A}}(\mathbf{X})\) , that is the set of elements of \(\hat{\mathrm{A}}(\mathbf{X})\) of the form \(a = \sum_{n\geqslant 1}a_n\) such that \(a_{n}\in \mathrm{L}^{n}(\mathrm{X})\) for all \(n\geqslant 0\) ; this is filtered Lie subalgebra of \(\hat{\mathrm{A}}(\mathbf{X})\) .

THEOREM 1. The restriction of the exponential mapping of \(\hat{\mathbf{A}}(\mathbf{X})\) to \(\hat{\mathbf{L}}(\mathbf{X})\) is a bijection of \(\hat{\mathbf{L}}(\mathbf{X})\) onto a closed subgroup of the Magnus group \(\Gamma(\mathbf{X})\) .

We write \(\mathrm{A}(\mathrm{X}) = \mathrm{A}\) , \(\mathrm{A}^{\mathrm{n}}(\mathrm{X}) = \mathrm{A}^{\mathrm{n}}\) , \(\hat{\mathrm{A}}(\mathrm{X}) = \hat{\mathrm{A}}\) , \(\mathrm{L}^{\mathrm{n}}(\mathrm{X}) = \mathrm{L}^{\mathrm{n}}\) , \(\hat{\mathrm{L}}(\mathrm{X}) = \hat{\mathrm{L}}\) , \(\Gamma(\mathrm{X}) = \Gamma\) . Let \(\mathrm{B}\) be the algebra \(\mathrm{A} \otimes \mathrm{A}\) with the graduation of type \(\mathbf{N}\) defined by \(\mathrm{B}^{\mathrm{n}} = \sum_{i+j=n} \mathrm{A}^{i} \otimes \mathrm{A}^{j}\) . Let \(\hat{\mathrm{B}} = \prod_{n > 0} \mathrm{B}^{n}\) be the associated complete filtered algebra (Commutative Algebra, Chapter III, §2, no. 12, Example 1). The co-product \(c: \mathrm{A} \to \mathrm{A} \otimes \mathrm{A}\) defined in §3, no. 1, Corollary 1 to Theorem 1 is graded of degree 0 and hence extends by continuity to a homomorphism \(\hat c: \hat{\mathrm{A}} \to \hat{\mathrm{B}}\) given by

\[
\hat {c} \left(\sum_ {n \geqslant 0} a _ {n}\right) = \sum_ {n \geqslant 0} c (a _ {n}) \quad \text { for } a _ {n} \in \mathrm{A} ^ {n}.
\]

We also define continuous homomorphisms \(\delta'\) and \(\delta''\) of \(\hat{\mathbf{A}}\) into \(\hat{\mathbf{B}}\) by

\[
\delta^ {\prime} \left(\sum_ {n \geqslant 0} a _ {n}\right) = \sum_ {n \geqslant 0} (a _ {n} \otimes 1), \quad \delta^ {\prime \prime} \left(\sum_ {n \geqslant 0} a _ {n}\right) = \sum_ {n \geqslant 0} (1 \otimes a _ {n}) \quad \text { for } a _ {n} \in \mathrm{A} ^ {n}.
\]

By Corollary 2 to Theorem 1 of §3, no. 1, \(\mathbf{L}^n\) is the set of \(a_{n}\in \mathbf{A}^{n}\) such that \(c(a_{n}) = a_{n}\otimes 1 + 1\otimes a_{n}\) . It follows that \(\hat{\mathbf{L}}\) is the set of \(a\in \hat{\mathbf{A}}\) such that

\[
\hat {c} (a) = \delta^ {\prime} (a) + \delta^ {\prime \prime} (a).\tag{3}
\]

Let \(\Delta\) be the set of \(b\in \hat{\mathbf{A}}\) of constant term equal to 1 and satisfying the relation

\[
\hat {c} (b) = \delta^ {\prime} (b). \delta^ {\prime \prime} (b),\tag{4}
\]

in other words, the set of \(b = \sum_{n\geqslant 0}b_n\) such that \(b_{n}\in \mathbf{A}^{n}\) for all \(n\geqslant 0\) , \(b_0 = 1\) and \(c(b_{n}) = \sum_{t + j = n}b_{t}\otimes b_{j}\) for \(n\geqslant 0\) . The latter characterization shows that \(\Delta\) is a closed subset of \(\Gamma\) ; as \(\hat c,\delta^{\prime}\) and \(\delta ''\) are ring homomorphisms and every element of \(\delta'(\hat{\mathbf A})\) commutes with every element of \(\delta''(\hat{\mathbf A})\) , the restrictions to \(\Gamma\) of the mappings \(c\) and \(\delta ' \delta''\) are group homomorphisms and \(\Delta\) is a subgroup of \(\Gamma\) .

By Proposition 1 of no. 1, the exponential mapping of \(\hat{\mathbf{A}}\) is a bijection of the set \(\hat{\mathbf{A}}^{+}\) of elements of \(\hat{\mathbf{A}}\) with no constant term onto \(\Gamma\) . Let \(a \in \hat{\mathbf{A}}^{+}\) and \(b = \exp a\) . As \(\hat c\) is a continuous ring homomorphism,

\[
\hat {c} (b) = \hat {c} \left(\sum_ {n \geqslant 0} a ^ {n} / n!\right) = \sum_ {n \geqslant 0} \hat {c} (a) ^ {n} / n! = \exp \hat {c} (a).
\]

The relations

\[
\delta^ {\prime} (b) = \exp \delta^ {\prime} (a), \quad \delta^ {\prime \prime} (b) = \exp \delta^ {\prime \prime} (a)
\]

are proved similarly and, as \(\delta'(a)\) commutes with \(\delta''(a)\) (no. 1, Remark 1),

\[
\delta^ {\prime} (b) \delta^ {\prime \prime} (b) = \exp (\delta^ {\prime} (a) + \delta^ {\prime \prime} (a)).
\]

Therefore a satisfies (3) if and only if b satisfies (4), which proves the theorem. Remark. The above proof shows that \(\exp(\hat{\mathbf{L}})\) is the subgroup \(\Delta\) of \(\Gamma\) consisting of the b satisfying (4).

Hence the group law of \(\Delta\) can be transported by the exponential mapping to \(\hat{\mathbf{L}}\) . In other words, \(\hat{\mathbf{L}}\) is a complete topological group with the law of composition \((a, b) \mapsto a \uplus b\) given by

\[
a \uplus b = \log (\exp a. \exp b).
\]

The topological group thus obtained is called the Hausdorff group (derived from X relative to K).

Let \(g\) be the homomorphism of the free group \(\mathbf{F} = \mathrm{F}(\mathbf{X})\) into \(\Gamma\) such that \(g(x) = \exp x\) for \(x \in \mathbf{X}\) . As \(\exp x - 1 - x = \sum_{n \geqslant 2} x^n / n!\) is of order \(\geqslant 2\) , \(g\) is injective by Theorem 1 of §5, no. 3. Therefore the mapping \(\log \circ g\) is an injective homomorphism of \(F\) into the Hausdorff group which extends the canonical injection \(X \to L\) .

For every integer \(m \geqslant 1\) we denote by \(\hat{\mathbf{L}}_m\) the set of elements of order \(\geqslant m\) in \(\hat{\mathbf{L}}\) and by \(\Gamma_m\) the set of \(u \in \Gamma\) such that \(u - 1\) is of order \(\geqslant m\) . Then \(\hat{\mathbf{L}}_m = \exp^{-1}(\Gamma_m)\) by Remark 2 of no. 1; as \((\Gamma_m)_{m \geqslant 1}\) is an integral central filtration on \(\Gamma\) (§4, no. 5 Proposition 2), \((\hat{\mathbf{L}}_m)_{m \geqslant 1}\) is an integral central filtration on the group \(\hat{\mathbf{L}}\) .

\subsection{LIE FORMAL POWER SERIES}

Lemma 1. Let \(\mathfrak{g}\) be a filtered Lie algebra (§ 4, no. 1), \((\mathfrak{g}_{\alpha})_{\alpha \in \mathbf{R}}\) its filtration and let \(\alpha \in \mathbf{R}\) . Let \(\mathbf{P}\) be a homogeneous Lie polynomial of degree \(n\) in the indeterminates \((\mathbf{T}_i)_{i \in I}\) (§ 2, no. 4). Then \(\mathrm{P}((a_i)) \in \mathfrak{g}_{n\alpha}\) for every family \((a_i)_{i \in I}\) of elements of \(\mathfrak{g}_{\alpha}\) .

Every Lie polynomial of degree \(n \geqslant 2\) is a finite sum of terms of the form \([\mathbf{Q}, \mathbf{R}]\) where \(\mathbf{Q}\) and \(\mathbf{R}\) are of degree \(< n\) and the sum their degrees is equal to \(n\) (§ 2, no. 7, Proposition 7). The lemma follows by induction on \(n\) .

A Lie formal power series† (with coefficients in K) in the indeterminates \((\mathrm{T}_{i})_{i\in\mathbb{I}}\) is any element of the Lie algebra \(\hat{\mathrm{L}}((\mathrm{T}_{i})_{i\in\mathbb{I}})=\hat{\mathrm{L}}(\mathrm{I})\) . Such an element u can be written uniquely as the sum of a summable family \((u_{v})_{v\in\mathbf{N}^{(\mathrm{I})}}\) where \(u_{v}\in\mathrm{L}^{v}(\mathrm{I})\) .

Suppose that I is finite. Let g be a complete Hausdorff filtered Lie algebra such that \(g = \bigcup_{\alpha > 0} g_{\alpha}\) ; let \(t = (t_i)_{i \in I}\) be a family of elements of g.

PROPOSITION 2. The homomorphism \(f_{t} \colon \mathrm{L}(\mathrm{I}) \to \mathfrak{g}\) such that \(f_{t}(\mathrm{T}_{i}) = t_{i}\) (§ 2, no. 4) can be extended by continuity to one and only one continuous homomorphism \(\hat{f}_{t}\) of \(\hat{\mathrm{L}}(\mathrm{I})\) into \(\mathfrak{g}\) .

There exists \(\alpha > 0\) such that \(t_i \in \mathfrak{g}_\alpha\) for all \(i \in I\) ; hence \(f_t(L^\nu(I)) \subset \mathfrak{g}_{|\nu| \alpha}\) for all \(\nu\) (Lemma 1), which implies the continuity of \(f_t\) .

If \(u \in \widehat{\mathbf{L}}(\mathbf{I})\) , we write \(u((t_i)) = \widehat{f_t}(u)\) . In particular, taking \(\mathfrak{g} = \widehat{\mathbf{L}}(\mathbf{I})\) , \(u = u((\mathrm{T}_i))\) ; in the general case, \(u((t_i))\) is called the result of substituting the \(t_i\) for the \(T_i\) in the Lie formal power series \(u((\mathrm{T}_i))\) . If \(u = \sum_{v \in \mathbf{N}^{(\mathrm{I})}} u_v\) , where \(u_v \in \mathrm{L}^v(\mathrm{I})\) , the family \((u_v((t_i))_{v \in \mathbf{N}^{(\mathrm{I})}})\) is summable and

\[
u ((t _ {i})) = \sum_ {\nu \in \mathbf {N} ^ {\mathrm{I}}} u _ {\nu} ((t _ {i})).\tag{5}
\]

Let \(\sigma\) be a continuous homomorphism of \(\mathfrak{g}\) into a complete Hausdorff filtered Lie algebra \(\mathfrak{g}'\) such that \(\mathfrak{g}' = \bigcup_{\alpha > 0} \mathfrak{g}_{\alpha}'\) . For every finite family \(\boldsymbol{t} = (t_i)_{i \in I}\) of elements of \(\mathfrak{g}\) and all \(u \in \hat{\mathbf{L}}(\mathbf{I})\) ,

\[
\sigma (u ((t _ {i}))) = u ((\sigma (t _ {i}))),\tag{6}
\]

for the homomorphism \(\sigma \circ f_t\) is continuous and maps \(T_i\) to \(\sigma(t_i)\) for \(i \in I\) .

Let \(\boldsymbol{u} = (u_j)_{j \in \mathcal{J}}\) be a finite family of elements of \(\hat{\mathrm{L}}(\mathrm{I})\) and let \(v \in \hat{\mathrm{L}}(\mathrm{J})\) ; by substituting the \(u_j\) for the \(T_j\) in \(v\) , we obtain an element \(w = v((u_j)_{j \in \mathcal{J}})\) of \(\hat{\mathrm{L}}(\mathrm{I})\) denoted by \(v \circ \boldsymbol{u}\) . Then

\[
w ((t _ {i}) _ {i \in I}) = v ((u _ {j} ((t _ {i}) _ {i \in I})) _ {j \in J})\tag{7}
\]

for every finite family \(\boldsymbol{t} = (t_{i})_{i \in I}\) of elements of g, as is seen by operating with the continuous homomorphism \(\hat{f}_{t}\) on the equation \(w = v((u_{j})_{j \in J})\) .

Let \(u = \sum_{\nu \in \mathbb{N}^{\mathbf{I}}} u_{\nu} \in \hat{\mathrm{L}}(\mathrm{I})\) , where \(u_{\nu} \in \mathrm{L}^{\nu}(\mathrm{I})\) . The mapping \(\tilde{u}\colon (t_i)\mapsto u((t_i))\) of \(\mathfrak{g}^{\mathbf{I}}\) into \(\mathfrak{g}\) is continuous: for in each of the open sets \(\mathfrak{g}_{\alpha}\) with \(\alpha > 0\) the family of \(\tilde{u}_{\nu}\) is uniformly summable and it suffices to prove that each \(\tilde{u}_{\nu}\) is continuous, which is immediate by induction on \(|\nu|\) .

\subsection{THE HAUSDORFF SERIES}

Let \(\{U, V\}\) be a set with two elements.

DEFINITION 1. The element \(\mathbf{H} = \mathbf{U} \uplus \mathbf{V} = \log (\exp \mathbf{U} \cdot \exp \mathbf{V})\) (no. 2) of the Lie algebra \(\hat{\mathbf{L}}_{\mathbf{Q}}(\{\mathbf{U}, \mathbf{V}\})\) is called the Hausdorff series in the indeterminates \(\mathbf{U}\) and \(\mathbf{V}\) .

We denote by \(H_{n}\) (resp. \(H_{r,s}\) ) the homogeneous component of H of total degree n (resp. multidegree \((r, s)\) ). Then

\[
\mathrm{H} = \sum_ {n \geqslant 0} \mathrm{H} _ {n} = \sum_ {r, s \geqslant 0} \mathrm{H} _ {r, s}, \quad \mathrm{H} _ {n} = \sum_ {\substack {r + s = n \\ r, s \geqslant 0}} \mathrm{H} _ {r, s}.\tag{8}
\]

THEOREM 2. If \(r\) and \(s\) are two positive integers such that \(r + s \geqslant 1\) , then \(\mathrm{H}_{r,s} = \mathrm{H}_{r,s}' + \mathrm{H}_{r,s}''\) , where

\[
(r + s) \mathrm{H} _ {r, s} ^ {\prime} = \tag {9}
\]

\[
\sum_{m\geqslant 1}\frac{(-1)^{m - 1}}{m}\sum_{\substack{r_{1} + \dots +r_{m} = r\\ s_{1} + \dots +s_{m - 1} = s - 1\\ r_{1} + s_{1}\geqslant 1,\ldots ,r_{m - 1} + s_{m - 1}\geqslant 1}}\left(\left(\prod_{i = 1}^{m - 1}\frac{(\mathrm{adU})^{r_{i}}}{r_{i}!}\frac{(\mathrm{adV})^{s_{i}}}{s_{i}!}\right)\frac{(\mathrm{adU})^{r_{m}}}{r_{m}!}\right)(\mathrm{V})
\]

\[
(r + s) \mathrm{H} _ {r, s} ^ {\prime \prime} = \tag {10}
\]

\[
\sum_{m\geqslant 1}\frac{(-1)^{m - 1}}{m}\sum_{\substack{r_{1} + \dots +r_{m - 1} = r - 1\\ s_{1} + \dots +s_{m - 1} = s\\ r_{1} + s_{1}\geqslant 1,\dots ,r_{m - 1} + s_{m - 1}\geqslant 1}}\Bigg(\prod_{i = 1}^{m - 1}\frac{(\operatorname{adU})^{r_{i}}}{r_{i}!}\frac{(\operatorname{adV})^{s_{i}}}{s_{i}!}\Bigg)(U).
\]

In \(\hat{\mathbf{A}}_{\mathbf{q}}(\{\mathrm{U},\mathrm{V}\})\) , exp U .exp V = 1 + W, where W = \(\sum_{r+s\geqslant 1}\frac{\mathrm{U}^r}{r!}\frac{\mathrm{V}^s}{s!}\) whence H = \(\sum_{m\geqslant 1}(-1)^{m - 1}\mathrm{W}^m /m\) (no. 2), that is:

\[
\mathrm{H}_{r,s} = \sum_{m\geqslant 1}\frac{(-1)^{m - 1}}{m}\sum_{\substack{r_{1} + \dots +r_{m} = r\\ s_{1} + \dots +s_{m} = s\\ r_{1} + s_{1}\geqslant 1,\ldots ,r_{m} + s_{m}\geqslant 1}}\prod_{i = 1}^{m}\frac{\mathrm{U}^{r_{i}}}{r_{i}!}\frac{\mathrm{V}^{s_{i}}}{s_{i}!}.\tag{11}
\]

The linear mapping \(\mathbf{P}_n\) , defined by \(\mathbf{P}_n(x_1, \ldots, x_n) = \frac{1}{n} \left( \prod_{i=1}^{n-1} (\text{ad } x_i) \right)(x_n)\) for \(n \geqslant 1\) and \(x_1, \ldots, x_n\) in \(\{\mathrm{U}, \mathrm{V}\}\) , is a projector of \(\mathrm{A}_{\mathbf{Q}}^n(\{\mathrm{U}, \mathrm{V}\})\) onto \(\mathrm{L}_{\mathbf{Q}}^n(\{\mathrm{U}, \mathrm{V}\})\) (§ 3, no. 2, Corollary to Proposition 1); as \(\mathrm{H}_{r,s}\) belongs to \(\mathbf{L}_{\mathbf{Q}}^{r+s}(\{\mathrm{U}, \mathrm{V}\})\) , \(\mathrm{H}_{r,s} = \mathbf{P}_{r+s}(\mathrm{H}_{r,s})\) . Now

\[
\begin{array}{l} \mathrm{P} _ {r + s} \left(\prod_ {i = 1} ^ {m} \frac {\mathrm{U} ^ {r _ {i}} \mathrm{V} ^ {s _ {i}}}{r _ {i} ! s _ {i} !}\right) \\ = \frac {1}{r + s} \left(\left(\prod_ {i = 1} ^ {m - 1} \frac {(\operatorname{ad} \mathrm{U}) ^ {r _ {i}}}{r _ {i} !} \frac {(\operatorname{ad} \mathrm{V}) ^ {s _ {i}}}{s _ {i} !}\right) \frac {(\operatorname{ad} \mathrm{U}) ^ {r _ {m}}}{r _ {m} !} \frac {(\operatorname{ad} \mathrm{V}) ^ {s _ {m} - 1}}{s _ {m} !}\right) (\mathrm{V}) \end{array}\tag{12}
\]

when \(s_m \geqslant 1\) and

\[
\mathrm{P} _ {r + s} \left(\prod_ {i = 1} ^ {m} \frac {\mathrm{U} ^ {r _ {i}}}{r _ {i} !} \frac {\mathrm{V} ^ {s _ {i}}}{s _ {i} !}\right) = \frac {1}{r + s} \left(\left(\prod_ {i = 1} ^ {m - 1} \frac {(\mathrm{adU}) ^ {r _ {i}}}{r _ {i} !} \frac {(\mathrm{adV}) ^ {s _ {i}}}{s _ {i} !}\right) \frac {(\mathrm{adU}) ^ {r _ {m} - 1}}{r _ {m} !}\right) (\mathrm{U})\tag{13}
\]

when \(r_m \geqslant 1\) and \(s_m = 0\) . Moreover, obviously (ad \(t\) ) \(^{p-1}\) . \(t = 0\) if \(p \geqslant 2\) and (ad \(t\) ) \(^0\) . \(t = t\) . It follows that the two sides of (12) are zero when \(s_m \geqslant 2\) and those of (13) are zero when \(r_m \geqslant 2\) . The theorem then follows since \(H_{r,s}'\) is the sum of the terms of type (12) and \(H_{r,s}''\) is the sum of the terms of type (13).

Remarks. (1) We have defined (§ 3, no. 2, Remark) a projector Q of A(X) onto L(X) such that \(Q(a^{m}) = 0\) for \(a \in L(X)\) and \(m \geqslant 2\) and Q(1) = 0. Then H = Q(exp H) = Q(exp U.exp V), whence immediately

\[
\mathrm{H} _ {r, s} = \mathrm{Q} \left(\frac {\mathrm{U} ^ {r}}{r !} \frac {\mathrm{V} ^ {s}}{s !}\right) \quad \text { for } r + s \geqslant 1.\tag{14}
\]

\begin{enumerate}
\def\labelenumi{(\arabic{enumi})}
\setcounter{enumi}{1}
\tightlist
\item
  We have
\end{enumerate}

\[
\begin{array}{r l} \mathrm{H(U,V)} & \equiv \mathrm{U+V} + \frac {1}{2} [ \mathrm{U,V} ] + \frac {1}{12} [ \mathrm{U,[U,V]} ] \\ & \quad + \frac {1}{12} [ \mathrm{V,[V,U]} ] - \frac {1}{24} [ \mathrm{U,[V,[U,V]} ] ] \end{array} \tag {15}
\]

modulo \(\sum_{n\geqslant 5}\mathrm{L}^n (\{\mathrm{U},\mathrm{V}\})\)

\begin{enumerate}
\def\labelenumi{(\arabic{enumi})}
\setcounter{enumi}{2}
\tightlist
\item
  \(\mathrm{H}_{0,n} = \mathrm{H}_{n,0} = 0\) for every integer \(n\neq 1\) , whence
\end{enumerate}

\[
\mathrm{H} (\mathrm{U}, 0) = \mathrm{H} (0, \mathrm{U}) = \mathrm{U}.\tag{16}
\]

On the other hand, as \([U, -U] = 0\) ,

\[
\mathrm{H} (\mathrm{U}, - \mathrm{U}) = 0.\tag{17}
\]

\subsection{SUBSTITUTIONS IN THE HAUSDORFF SERIES}

As K is a field containing Q, the Hausdorff series can be considered as a Lie formal power series with coefficients in K. Therefore, if g is a complete Hausdorff filtered Lie algebra with \(g = \bigcup_{\alpha > 0} g_{\alpha}\) , then, for a, b in g, a and b can be substituted for U and V in H (cf.~no. 3 and §2, no. 5, Remark).

In particular, let A be a complete Hausdorff filtered unital associative algebra. We write \(m = \bigcup_{\alpha > 0} A_{\alpha}\) and \(m_{\alpha} = A_{\alpha} \cap m\) for \(\alpha \in R\) ; hence \(m_{\alpha} = A_{\alpha}\) for \(\alpha > 0\) and \(m_{\alpha} = m\) for \(\alpha \leqslant 0\) . With the bracket \([a, b] = ab - ba\) , m is a complete Hausdorff filtered Lie algebra, to which the above can be applied. With this notation, we have the following result which completes Proposition 1 of no. 1.

PROPOSITION 3. If \(a \in \mathfrak{m}\) , \(b \in \mathfrak{m}\) , then \(\exp \mathrm{H}(a, b) = \exp a \cdot \exp b\) .

Let \(a, b\) be in \(\mathfrak{m}\) ; there exists \(\alpha > 0\) such that \(a \in \mathrm{A}_{\alpha}\) and \(b \in \mathrm{A}_{\alpha}\) . Then there exists a continuous homomorphism \(\theta\) of the Magnus algebra \(\hat{\mathbf{A}}(\{\mathbf{U},\mathbf{V}\})\) into A mapping U to \(a\) and V to \(b\) ( \(\S 5\) , no. 1, Proposition 1).

The restriction of \(\theta\) to \(\hat{\mathbf{L}}(\{\mathrm{U},\mathrm{V}\})\) is a continuous homomorphism of Lie algebras of \(\mathbf{L}(\{\mathrm{U},\mathrm{V}\})\) into \(\mathfrak{m}\) which maps \(\mathbf{U}\) (resp. V) to \(a\) (resp. \(b\) ). By formula (6) of no. 3, therefore \(\theta (\mathrm{H}) = \mathrm{H}(a,b)\) . It then suffices to apply the continuous homomorphism \(\theta\) to the two sides of the relation

\[
\exp \mathrm{H} (\mathrm{U}, \mathrm{V}) = \exp \mathrm{U}. \exp \mathrm{V}
\]

taking account of Remark 3 of no. 1.

Remark 1. If \(a\) and \(b\) commute, then \(\mathrm{H}_{r,s}(a,b) = 0\) for \(r + s \geqslant 2\) , for every homogeneous Lie polynomial of degree \(\geqslant 2\) is zero at \((a,b)\) . Then \(\mathrm{H}(a,b) = a + b\) and Proposition 3 recovers the formula

\[
\exp (a + b) = \exp a. \exp b.
\]

PROPOSITION 4. Let \(\mathfrak{g}\) be a complete Hausdorff filtered Lie algebra such that \(\mathfrak{g} = \bigcup_{\alpha >0} \mathfrak{g}_{\alpha}\) . The mapping \((a, b) \mapsto \mathrm{H}(a, b)\) is a group law on \(\mathfrak{g}\) compatible with the topology on \(\mathfrak{g}\) under which \(0\) is the identity element and \(-a\) is the inverse of \(a\) for all \(a \in \mathfrak{g}\) .

The mapping \((a, b) \mapsto \mathrm{H}(a, b)\) of \(\mathfrak{g} \times \mathfrak{g}\) into \(\mathfrak{g}\) is continuous (no. 3); as the mapping \(a \mapsto -a\) is obviously continuous, it suffices to prove the relations

\begin{enumerate}
\def\labelenumi{(\arabic{enumi})}
\setcounter{enumi}{17}
\tightlist
\item
\end{enumerate}

\[
\mathrm{H} (\mathrm{H} (a, b), c) = \mathrm{H} (a, \mathrm{H} (b, c))\tag{18}
\]

\[
\mathrm{H} (a, - a) = 0\tag{19}
\]

\[
\mathrm{H} (a, 0) = \mathrm{H} (0, a) = a\tag{20}
\]

for \(a, b, c\) in \(\mathfrak{g}\) . By formula (7) of no. 3, it suffices to prove these formulae when \(a, b, c\) are three indeterminates and \(\mathfrak{g} = \hat{\mathbf{L}}\langle \{a, b, c\}\rangle\) . Now the restriction of the exponential mapping to \(\hat{\mathbf{L}}\langle \{a, b, c\}\rangle\) is an injection into the Magnus algebra \(\hat{\mathbf{A}}\langle \{a, b, c\}\rangle\) and by Proposition 3:

\[
\exp \mathrm{H} (\mathrm{H} (a, b), c) = \exp \mathrm{H} (a, b). \exp c = \exp a. \exp b. \exp c
\]

\[
\exp \mathrm{H} (a, \mathrm{H} (b, c)) = \exp a. \exp \mathrm{H} (b, c) = \exp a. \exp b. \exp c
\]

\[
\exp \mathrm{H} (a, - a) = \exp a. \exp (- a) = \exp (a - a) = \exp 0
\]

\[
\exp \mathrm{H} (a, 0) = \exp a. \exp 0 = \exp a
\]

\[
\exp \mathrm{H} (0, a) = \exp 0. \exp a = \exp a.
\]

This establishes relations (18) to (20).

Remarks. (2) Take g to be the Lie algebra \(\hat{\mathbf{L}}(\mathbf{X})\) . The group law introduced in the above proposition coincides with the law defined in no. 2. In other words,

\[
a \uplus b = \mathrm{H} (a, b) \quad \text { for   } a, b \text {   in   } \hat {\mathrm{L}} (\mathrm{X});\tag{21}
\]

thus the Hausdorff group law is given by the Hausdorff series.

\begin{enumerate}
\def\labelenumi{(\arabic{enumi})}
\setcounter{enumi}{2}
\tightlist
\item
  Let \(\mathfrak{g}\) be a Lie algebra with the integral filtration \((\mathcal{C}^m\mathfrak{g})\) defined by the lower central series. Suppose that there exists \(m \geqslant 1\) such that \(\mathcal{C}^m\mathfrak{g} = \{0\}\) . With the topology derived from the filtration \((\mathcal{C}^m\mathfrak{g})_{n \geqslant 1}\) , the Lie algebra \(\mathfrak{g}\) is Hausdorff, complete and even discrete. Then \(\mathrm{P}(a_1, \ldots, a_r) = 0\) for \(a_1, \ldots, a_r\) in \(\mathfrak{g}\) and for every homogeneous Lie polynomial \(\mathrm{P}\) of degree \(\geqslant m\) ; in particular, \(\mathrm{H}_{r,s}(a, b) = 0\) for \(r + s \geqslant m\) and the series \(\mathrm{H}(a, b) = \sum_{r,s} \mathrm{H}_{r,s}(a, b)\) has only a finite number of non-zero terms. The group law \((a, b) \mapsto \mathrm{H}(a, b)\) on \(\mathfrak{g}\) is then a polynomial mapping (§ 2, no. 4).
\end{enumerate}

PROPOSITION 5. Let \(\mathbf{K}_{r,s}\) be the component of \(\mathrm{H}(\mathrm{U} + \mathrm{V}, -\mathrm{U})\) of multidegree \((r,s)\) . Then

\[
\mathrm{K} _ {n, 1} (\mathrm{U}, \mathrm{V}) = \frac {1}{(n + 1) !} (\operatorname{ad} \mathrm{U}) ^ {n} (\mathrm{V}) \quad for n \geqslant 0.
\]

We write \(\mathrm{K}(\mathrm{U},\mathrm{V}) = \mathrm{H}(\mathrm{U} + \mathrm{V}, - \mathrm{U}),\mathrm{K}_1(\mathrm{U},\mathrm{V}) = \sum_{n\geqslant 0}\mathrm{K}_{n,1}(\mathrm{U},\mathrm{V}).\) We denote by L (resp. R) left (resp. right) multiplication by U on \(\hat{\mathbf{A}} (\{\mathbf{U},\mathbf{V}\})\)

We can write

\[
\begin{array}{r l} e ^ {\mathrm{U}} \mathrm{V} e ^ {- \mathrm{U}} & = \sum_ {p, q} \frac {\mathrm{U} ^ {p}}{p !} \mathrm{V} \frac {(- \mathrm{U}) ^ {q}}{q !} \\ & = \sum_ {n \geqslant 0} \frac {1}{n !} \left(\sum_ {p + q = n} \frac {n !}{p ! q !} (\mathrm{L} ^ {p} (- \mathrm{R}) ^ {q}). \mathrm{V}\right) \\ & = \sum_ {n \geqslant 0} \frac {1}{n !} (\mathrm{L} - \mathrm{R}) ^ {n}. \mathrm{V} \end{array}
\]

and therefore

\[
e ^ {\mathrm{U}} \mathrm{V} e ^ {- \mathrm{U}} = \sum_ {n \geqslant 0} \frac {1}{n !} (\mathrm{adU}) ^ {n} \mathrm{V}.\tag{22}
\]

We now calculate modulo the ideal \(\sum_{m\geqslant 0}\sum_{n\geqslant 2}\mathrm{A}^{m,n}(\{\mathrm{U},\mathrm{V}\})\) of \(\mathrm{A}(\{\mathrm{U},\mathrm{V}\})\) . For all \(n\geqslant 1\) ,

\[
(\mathrm{U} + \mathrm{V}) ^ {n} \equiv \mathrm{U} ^ {n} + \sum_ {i = 1} ^ {n - 1} \mathrm{U} ^ {i} \mathrm{V} \mathrm{U} ^ {n - 1 - i}
\]

whence

\[
\begin{array}{r l} (\mathrm{adU}) (\mathrm{U} + \mathrm{V}) ^ {n} & \equiv ((\mathrm{L} - \mathrm{R}) \sum_ {i = 1} ^ {n - 1} \mathrm{L} ^ {i} \mathrm{R} ^ {n - i}). \mathrm{V} \\ & \equiv (\mathrm{L} ^ {n} - \mathrm{R} ^ {n}). \mathrm{V} \\ & \equiv \mathrm{U} ^ {n} \mathrm{V} - \mathrm{VU} ^ {n}. \end{array}
\]

Therefore

\[
(\mathrm{ad} \mathrm{U}). \mathrm{e} ^ {\mathrm{U} + \mathrm{V}} \equiv e ^ {\mathrm{U}} \mathrm{V} - \mathrm{V} e ^ {\mathrm{U}}\tag{23}
\]

summing over n.

On the other hand, \(\mathrm{K}_1(\mathrm{U},\mathrm{V})\equiv \mathrm{K}(\mathrm{U},\mathrm{V})\) and \(e^{\mathrm{K}_1(\mathrm{U},\mathrm{V})}\equiv 1 + \mathrm{K}_1(\mathrm{U},\mathrm{V})\) and hence

\[
\mathrm{K} _ {1} \equiv e ^ {\mathrm{K}} - 1 \equiv e ^ {\mathrm{U} + \mathrm{V}} e ^ {- \mathrm{U}} - 1
\]

by Proposition 3. We deduce that

\[
\begin{array}{r l r} (\mathrm{adU}) \mathrm{K} _ {1} & \equiv \mathrm{Ue} ^ {\mathrm{U} + \mathrm{V}} e ^ {- \mathrm{U}} - e ^ {\mathrm{U} + \mathrm{V}} e ^ {- \mathrm{U}} \mathrm{U} \equiv (\mathrm{Ue} ^ {\mathrm{U} + \mathrm{V}} - e ^ {\mathrm{U} + \mathrm{V}} \mathrm{U}) e ^ {- \mathrm{U}} \\ & \equiv (e ^ {\mathrm{U}} \mathrm{V} - \mathrm{Ve} ^ {\mathrm{U}}) e ^ {- \mathrm{U}} & \text {by (23)} \\ & \equiv e ^ {\mathrm{U}} \mathrm{Ve} ^ {- \mathrm{U}} - \mathrm{V} \\ & \equiv \sum_ {n \geqslant 1} \frac {1}{n !} (\mathrm{adU}) ^ {n} \mathrm{V} & \text {by (22)} \\ & \equiv (\mathrm{adU}) \Bigl (\sum_ {n \geqslant 0} \frac {(\mathrm{adU}) ^ {n}}{(n + 1) !} \mathrm{V} \Bigr). \end{array}
\]

It then suffices to apply the Remark of § 2, no. 11.

\section{CONVERGENCE OF THE HAUSDORFF SERIES (REAL OR COMPLEX CASE)}

In this paragraph we assume that K is one of the fields R or C with its usual absolute value. Recall that a normable algebra over K is a (not necessarily associative) algebra over K with a topology T with the following properties:

\begin{enumerate}
\def\labelenumi{(\arabic{enumi})}
\item
  \(\mathcal{T}\) can be defined by a norm:
\item
  the mapping \((x,y)\mapsto xy\) of \(\mathbf{A}\times \mathbf{A}\) into \(\mathbf{A}\) is continuous.
\end{enumerate}

A normed algebra over K is an algebra A over K with a norm such that \(\|xy\|\leqslant\|x\|\|y\|\) for all x,y in A.

We denote by g a complete normable Lie algebra over K. We choose a norm on g and a number M > 0 such that

\[
\| [ x, y ] \| \leqslant \mathrm{M} \| x \| \| y \| \quad \text { for } x, y \text { in } \mathfrak {g}.\tag{1}
\]

\subsection{CONTINUOUS-POLYNOMIALS WITH VALUES IN g}

Let I be a finite set and let \(\mathbf{P}(\mathfrak{g}^{\mathrm{I}};\mathfrak{g})\) (resp. \(\hat{\mathbf{P}} (\mathfrak{g}^{\mathrm{I}};\mathfrak{g}))\) be the vector space of continuous-polynomials (resp. formal power series with continuous components) on \(\mathbf{g}^{\mathrm{I}}\) with values in \(\mathfrak{g}\) . Recall (Differentiable and Analytic Manifolds, R, Appendix) that \(\mathbf{P}(\mathfrak{g}^{\mathrm{I}};\mathfrak{g})\) has a graduation of type \(\mathbf{N}^{\mathrm{I}}\) and that \(\hat{\mathbf{P}} (\mathfrak{g}^{\mathrm{I}};\mathfrak{g})\) is identified with the completion of the vector space \(\mathbf{P} (\mathfrak{g}^{\mathrm{I}};\mathfrak{g})\) with the topology defined by the filtration associated with the graduation of \(\mathbf{P}(\mathfrak{g}^{\mathrm{I}};\mathfrak{g})\) . Moreover, \(\mathbf{P}(\mathfrak{g}^{\mathrm{I}};\mathfrak{g})\) is a graded Lie algebra with the bracket defined by \([f,g](\pmb {x}) = [f(\pmb {x}),g(\pmb {x})]\) for \(f, g\) in \(\mathrm{P}(\mathfrak{g}^{\mathrm{I}}; \mathfrak{g}), x \in \mathfrak{g}^{\mathrm{I}}\) ; this Lie algebra structure can be extended by continuity to \(\hat{\mathrm{P}}(\mathfrak{g}^{\mathrm{I}}; \mathfrak{g})\) and makes it into a complete Hausdorff filtered Lie algebra.

By proposition 2 of § 6, no. 3, there exists one and only one continuous Lie algebra homomorphism \(\phi_{\mathrm{I}}:u\mapsto \tilde{u}\) of \(\hat{\mathbf{L}} (\mathbf{I})\) into \(\hat{\mathbf{P}} (\mathbf{g}^{\mathrm{I}};\mathfrak{g})\) mapping the indeterminate of index \(i\) to \(\mathsf{pr}_{\mathrm{i}}\) for all \(i\in \mathbb{I}\) , since \(\mathsf{pr}_{\mathrm{i}}\in \mathsf{P}(\mathfrak{g}^{\mathrm{I}};\mathfrak{g})\) . It follows that \(\tilde{u}\in \mathsf{P}(\mathfrak{g}^{\mathrm{I}};\mathfrak{g})\) for \(u\in \mathbf{L}(\mathbf{I})\) ; more precisely, when \(u\in \mathbf{L}(\mathbf{I})\) , \(\tilde{u}\) is just the polynomial mapping \((t_i)\mapsto u((t_i))\) of § 2, no. 4. On the other hand, clearly \(\phi_{\mathrm{I}}\) is compatible with the multigraduations of \(\mathbf{L}(\mathbf{I})\) and \(\mathbf{P}(\mathfrak{g}^{\mathrm{I}};\mathfrak{g})\) . If \(u = \sum_{\nu \in \mathbf{N}^{\mathrm{I}}}u_{\nu}\) , where \(u_{\nu}\in \mathbf{L}^{\nu}(\mathbf{I})\) for \(\nu \in \mathbf{N}^{\mathrm{I}}\) , then

\[
\tilde {u} = \sum_ {\nu \in \mathbf {N} ^ {\mathrm{I}}}  \tilde {u} _ {\nu}, \quad \text { where } \tilde {u} _ {\nu} \in \mathrm{P} _ {\nu} (\mathfrak {g} ^ {\mathrm{I}}; \mathfrak {g}).
\]

Let \(\pmb{u} = (u_j)_{j\in J}\) be a finite family of elements of \(\hat{\mathbf{L}} (\mathbf{I})\) , let \(v\in \hat{\mathbf{L}} (\mathbf{J})\) and let \(w = v\circ \pmb {u}\) (§ 6, no. 3). We write \(\tilde{\pmb{u}} = (\tilde{u}_j)_{j\in \mathbb{J}}\in \mathfrak{J}\) . Then

\[
\tilde {v} \circ \tilde {\boldsymbol {u}} = (v \circ \boldsymbol {u}) ^ {\sim}.\tag{2}
\]

This follows by extending by continuity formula (7) of §6, no. 3 and from Differentiable and Analytic Manifolds, R, Appendix, no. 6.

\subsection{GROUP GERM DEFINED BY A COMPLETE NORMED LIE ALGEBRA}

Let \(H = \sum_{r,s \geqslant 0} H_{r,s} \in \hat{L}(U, V)\) be the Hausdorff series (§ 6, no. 4, Definition 1). We shall show that the corresponding formal power series

\[
\tilde {\mathrm{H}} = \sum_ {r, s \geqslant 0} \tilde {\mathrm{H}} _ {r, s} \in \hat {\mathrm{P}} (\mathfrak {g} \times \mathfrak {g}, \mathfrak {g})\tag{3}
\]

is convergent (Differentiable and Analytic Manifolds, R, 3.1.1).

We introduce the following formal power series \(\eta\in\mathbf{Q}[[U,V]]\)

\begin{enumerate}
\def\labelenumi{(\arabic{enumi})}
\setcounter{enumi}{3}
\tightlist
\item
\item
\end{enumerate}

\[
\begin{array}{l}\eta (\mathrm{U},\mathrm{V}) = -\log (2 - \exp (\mathrm{U} + \mathrm{V}))\\ \qquad = \sum_{m\geqslant 1}\frac{1}{m} (\exp (\mathrm{U} + \mathrm{V}) - 1)^{m}\\ \qquad = \sum_{m\geqslant 1}\frac{1}{m}\sum_{\substack{r_{1},\ldots ,r_{m}\\ s_{1},\ldots ,s_{m}\\ r_{i} + s_{i}\geqslant 1}}\frac{\mathrm{U}^{r_{1}}}{r_{1}!}\frac{\mathrm{V}^{s_{1}}}{s_{1}!}\frac{\mathrm{U}^{r_{2}}}{r_{2}!}\dots \frac{\mathrm{V}^{s_{m}}}{s_{m}!}. \end{array}\tag{6}
\]

Hence

\[
\eta (\mathrm{U}, \mathrm{V}) = \sum_ {r, s \geqslant 0} \eta_ {r, s} \mathrm{U} ^ {r} \mathrm{V} ^ {s},\tag{7}
\]

where

\[
\eta_{r,s} = \sum_{m\geqslant 1}\frac{1}{m}\sum_{\substack{r_{1} + \dots +r_{m} = r\\ s_{1} + \dots +s_{m} = s\\ r_{i} + s_{i}\geqslant 1}}\frac{1}{r_{1}!\ldots r_{m}!s_{1}!\ldots s_{m}!}.\tag{8}
\]

Now let u and v be two positive real numbers such that \(u + v < \log 2\) ; then \(0 \leqslant \exp(u + v) - 1 < 1\) ; the series derived from (5) and (6) by substituting u for U and v for V are convergent and the above calculations imply that

\[
\sum_ {r, s \geqslant 0} \eta_ {r, s} u ^ {r} v ^ {s} = - \log (2 - \exp (u + v)) <   + \infty .\tag{9}
\]

Let \(r, s \geqslant 0\) and let \(\| \tilde{\mathrm{H}}_{r,s} \|\) be the norm of the continuous-polynomial \(\tilde{\mathrm{H}}_{r,s}\) (Differentiable and Analytic Manifolds, R, Appendix, no. 2).

Lemma 1.

\[
\| \widetilde {\mathrm{H}} _ {r, s} \| \leqslant \mathrm{M} ^ {r + s - 1} \eta_ {r, s}.
\]

Let \(r_i, s_i\) be in \(\mathbf{N}\) for \(1 \leqslant i \leqslant m\) , with \(s_m = 1\) ; we write \(r = \sum_{i} r_i, s = \sum_{i} s_i\) and consider the following element of \(L(\{U, V\})\) :

\[
Z = \left(\left(\sum_ {i = 1} ^ {m - 1} (\operatorname{ad} U) ^ {r _ {i}} (\operatorname{ad} V) ^ {s _ {i}}\right) (\operatorname{ad} U) ^ {r _ {m}}\right) (V).
\]

Then \(\tilde{Z} = f \circ p\) , where \(f\) is the following \((r + s)\) -linear mapping of \(\mathfrak{g}^{r + s}\) into \(\mathfrak{g}\) :

\[
\begin{array}{l} \left(x _ {1}, \dots , x _ {r}, y _ {1}, \dots , y _ {s}\right) \mapsto \\ \left(\operatorname{ad} \left(x _ {1}\right) \circ \dots \circ \operatorname{ad} \left(x _ {r _ {1}}\right) \circ \operatorname{ad} \left(y _ {1}\right) \circ \dots \circ \operatorname{ad} \left(y _ {s _ {1}}\right) \circ \operatorname{ad} \left(x _ {r _ {1} + 1}\right) \circ \dots \circ \operatorname{ad} \left(x _ {r}\right)\right) \left(y _ {s}\right) \end{array}
\]

and where \(p\) is the following mapping of \(\mathfrak{g}^2\) into \(\mathfrak{g}^{r+s}\) :

\[
(x, y) \mapsto \underbrace {(x , \dots , x} _ {r}, \underbrace {y , \dots , y)} _ {s};
\]

hence \(\| \tilde{Z}\| \leqslant \| f\| \leqslant M^{r + s - 1}\) (Differentiable and Analytic Manifolds, R, Appendix). Applying these inequalities to the various terms on the right hand side of formula (9) of §6, no. 4, we obtain:

\[
\| (\mathrm{H} _ {r, s} ^ {\prime}) ^ {\sim} \| \leqslant \frac {\mathrm{M} ^ {r + s - 1}}{r + s} \sum_ {m \geqslant 1} \frac {1}{m} \sum_ {\substack {r _ {1} + \dots + r _ {m} = r \\ s _ {1} + \dots + s _ {m - 1} = s - 1 \\ r _ {1} + s _ {1} \geqslant 1, \dots , r _ {m - 1} + s _ {m - 1} \geqslant 1}} \frac {1}{r _ {1} ! \dots r _ {m} ! s _ {1} ! \dots s _ {m - 1} !}.\tag{10}
\]

A similar argument gives

\[
\begin{array}{l}\| (\mathbf{H}_{r,s}^{\prime \prime})\sim \| \\ \leqslant \frac{\mathbf{M}^{r + s - 1}}{r + s}\sum_{m\geqslant 1}\frac{1}{m}\sum_{\substack{r_{1} + \dots +r_{m - 1} = r - 1\\ s_{1} + \dots +s_{m - 1} = s\\ r_{1} + s_{1}\geqslant 1,\ldots ,r_{m - 1} + s_{m - 1}\geqslant 1}}\frac{1}{r_{1}!\ldots r_{m - 1}!s_{1}!\ldots s_{m - 1}!} \end{array}\tag{11}
\]

whence, by (8)

\[
\left\| \tilde {\mathrm{H}} _ {r, s} \right\| \leqslant   \eta_ {r, s} \frac {\mathrm{M} ^ {r + s - 1}}{r + s} \leqslant \eta_ {r, s} \mathrm{M} ^ {r + s - 1},
\]

which proves the lemma.

PROPOSITION 1. The formal power series \(\tilde{\mathbf{H}}\) is a convergent series (Differentiable and Analytic Manifolds, R, 3.1.1); its domain of absolute convergence (Differentiable and Analytic Manifolds, R, 3.1.4) contains the open set

\[
\Omega = \left\{(x, y) \in \mathfrak {g} \times \mathfrak {g} | \| x \| + \| y \| <   \frac {1}{M} \log 2 \right\}.
\]

Let \(u, v\) be two real numbers \(> 0\) such that \(u + v < \frac{1}{M} \log 2\) ; then (Lemma 1)

\[
\begin{array}{r l} \mathrm{M} \sum_ {r, s \geqslant 0} \| \tilde {\mathrm{H}} _ {r, s} \| u ^ {r} v ^ {s} \\ & \leqslant \sum_ {r, s \geqslant 0} \eta_ {r, s} \mathrm{M} ^ {r + s} u ^ {r} v ^ {s} = - \log (2 - \exp \mathrm{M} (u + v)) <   + \infty \end{array} \tag {12}
\]

by (9).

Let \(h: \Omega \to \mathfrak{g}\) denote the analytic function (Differentiable and Analytic Manifolds, R, 3.2.9) defined by \(\tilde{\mathrm{H}}\) , that is by the formula

\[
h (x, y) = \sum_ {r, s \geqslant 0} \tilde {\mathrm{H}} _ {r, s} (x, y) = \sum_ {r, s \geqslant 0} \mathrm{H} _ {r, s} (x, y) \quad \text { for } (x, y) \in \Omega .\tag{13}
\]

This function is called the Hausdorff function of g relative to M (or simply the Hausdorff function of g if no confusion can arise). Note that \(\mathrm{H}_{r,s}(\mathrm{U}, -\mathrm{U}) = 0\) if \(r + s \geqslant 2\) and hence

\[
h (x, - x) = 0 \quad \text { for } \| x \| <   \frac {1}{2 M} \log 2.\tag{14}
\]

Similarly

\[
h (0, x) = h (x, 0) = x \quad \text { for } \| x \| <   \frac {1}{\mathrm{M}} \log 2.\tag{15}
\]

PROPOSITION 2. Let

\[
\Omega^ {\prime} = \left\{(x, y, z) \in \mathfrak {g} \times \mathfrak {g} \times \mathfrak {g} | \| x \| + \| y \| + \| z \| <   \frac {1}{M} \log \frac {3}{2} \right\}.
\]

If \((x,y,z)\in \Omega^{\prime}\) , then

\[
(x, y) \in \Omega , \quad (h (x, y), z) \in \Omega , \quad (y, z) \in \Omega , \quad (x, h (y, z)) \in \Omega\tag{16}
\]

and

\[
h (h (x, y), z) = h (x, h (y, z)).\tag{17}
\]

Let \((x,y,z)\in \Omega^{\prime}\) ; clearly \((x,y)\in \Omega\) and \((y,z)\in \Omega\) . Moreover:

\[
\| h (x, y) \| \leqslant \sum_ {r, s} \| \tilde {\mathrm{H}} _ {r, s} \| \| x \| ^ {r} \| y \| ^ {s},
\]

and hence by (13)

\[
\| h (x, y) \| \leqslant - \frac {1}{M} \log (2 - \exp M (\| x \| + \| y \|)).
\]

Now \(\mathrm{M}(\|x\| + \|y\|) < \log \frac{3}{2} - \mathrm{M}\|z\|\) ; we write \(u = \exp(\mathrm{M}\|z\|)\) ; then \(1 \leqslant u \leqslant \frac{3}{2}\) and

\[
\begin{array}{r l} \mathbf {M} (\| h (x, y) \| + \| z \|) <   - \log (2 - \exp (\log \frac {3}{2} - \mathbf {M} \| z \|)) + \mathbf {M} \| z \| \\ = - \log \left(2 - \frac {3}{2 u}\right) + \log u = \log \frac {2 u ^ {2}}{4 u - 3} \\ = \log \left(2 + \frac {2 (u - 1) (u - 3)}{4 u - 3}\right) \leqslant \log 2. \end{array}
\]

We see similarly that \((x, h(y, z)) \in \Omega\) .

We now prove (17). In the Lie algebra \(\hat{\mathbf{L}}(\{\mathbf{U},\mathbf{V},\mathbf{W}\})\)

\[
\mathrm{H} (\mathrm{H} (\mathrm{U}, \mathrm{V}), \mathrm{W}) = \mathrm{H} (\mathrm{U}, \mathrm{H} (\mathrm{V}, \mathrm{W}))
\]

by Proposition 4 of §6, no. 5. By no. 1, formula (2), we therefore have in \(\hat{\mathbf{P}}(\mathfrak{g} \times \mathfrak{g} \times \mathfrak{g}, \mathfrak{g})\) the relation

\[
\tilde {\mathrm{H}} \circ (\tilde {\mathrm{H}} \times \mathrm{Id} _ {\mathfrak {g}}) = \tilde {\mathrm{H}} \circ (\mathrm{Id} _ {\mathfrak {g}} \times \tilde {\mathrm{H}}).
\]

By Differentiable and Analytic Manifolds, R, 3.1.9, there exists a number \(\varepsilon > 0\) such that formula (17) is true when \(\|x\|\) , \(\|y\|\) and \(\|z\|\) are \(\leqslant \varepsilon\) . But the functions \((x,y,z)\mapsto h(h(x,y),z)\) and \((x,y,z)\mapsto h(x,h(y,z))\) are analytic functions on \(\Omega'\) with values in \(\mathfrak{g}\) (Differentiable and Analytic Manifolds, R, 3.2.7). As \(\Omega'\) is connected and they coincide in a neighbourhood of 0, they are equal (Differentiable and Analytic Manifolds, R, 3.2.5).

The above results imply:

Let \(\alpha\) be a real number such that \(0 < \alpha \leqslant \frac{1}{3M} \log \frac{3}{2}\) . Let

\[
\mathrm{G} = \{x \in \mathfrak {g} \mid \| x \| <   \alpha \},
\]

\(\Theta = \{(x,y) \in G \times G \mid h(x,y) \in G\}\) and \(m: \Theta \to G\) be the restriction of \(h\) to \(\Theta\) . Then:

\begin{enumerate}
\def\labelenumi{(\arabic{enumi})}
\item
  \(\Theta\) is open in \(\mathbf{G} \times \mathbf{G}\) and \(m\) is analytic.
\item
  \(x \in G\) implies \((0, x) \in \Theta\) , \((x, 0) \in \Theta\) and \(m(0, x) = m(x, 0) = x\) .
\item
  \(x \in \mathbf{G}\) implies \(-x \in \mathbf{G}, (x, -x) \in \Theta, (-x, x) \in \Theta\) and
\end{enumerate}

\[
m (x, - x) = m (- x, x) = 0.
\]

\begin{enumerate}
\def\labelenumi{(\arabic{enumi})}
\setcounter{enumi}{3}
\tightlist
\item
  Let \(x, y, z\) be elements of \(G\) such that \((x, y) \in \Theta\) , \((m(x, y), z) \in \Theta\) , \((y, z) \in \Theta\) and \((x, m(y, z)) \in \Theta\) . Then \(m(m(x, y), z) = m(x, m(y, z))\) .
\end{enumerate}

*In other words (Chapter III, §1), if we write \(-x = \sigma(x)\) , the quadruple \((G, 0, \sigma, m)\) is a Lie group germ over \(K_*\)

\subsection{EXPONENTIAL IN COMPLETE NORMED ASSOCIATIVE ALGEBRAS}

In this no. we denote by A a complete normed unital associative algebra (General Topology, Chapter IX, § 3, no. 7). Then \(\| x.y\| \leqslant \| x\|\) . \(\| y\|\) for \(x,y\) in A.

Let I be a finite set and let \(\hat{\mathrm{P}} (\mathrm{A}^{\mathrm{I}};\mathrm{A})\) be the vector space of formal power series with continuous components on \(\mathbf{A}^{\mathrm{I}}\) with values in A (Differentiable and Analytic Manifolds, R, Appendix, no. 5) with the algebra structure obtained by writing

\[
f. g = m \circ (f, g) \quad \text { for } f, g \text { in } \hat {\mathrm{P}} (\mathrm{A} ^ {\mathrm{I}}; \mathrm{A}),
\]

where \(m: \mathbf{A} \times \mathbf{A} \to \mathbf{A}\) denotes multiplication on A. Arguing as in no. 1 and using Proposition 1 of § 5, no. 1, we define a continuous homomorphism of unital algebras \(u \mapsto \tilde{u}\) of \(\hat{\mathrm{A}}(\mathrm{I})\) into \(\hat{\mathrm{P}}(\mathrm{A}^{\mathrm{I}}; \mathrm{A})\) mapping the indeterminate of index \(i\) to \(\mathbf{pr}_{\mathrm{i}}\) ; this homomorphism extends the Lie algebra homomorphism of \(\hat{\mathrm{L}}(\mathrm{I})\) into \(\hat{\mathrm{P}}(\mathrm{A}^{\mathrm{I}}; \mathrm{A})\) defined in no. 1. If \(u = \sum_{\nu} u_{\nu}\) with \(u_{\nu} \in \mathrm{A}^{\nu}(\mathrm{I})\) for \(\nu \in \mathbf{N}^{\mathrm{I}}\) , then \(\tilde{u} = \sum_{\nu} \tilde{u}_{\nu}\) , where \(\tilde{u}_{\nu}\) is the polynomial mapping \((t_{i})_{i \in I} \mapsto u_{\nu}((t_{i}))\) .

Let \(\boldsymbol{u} = (u_{j})_{j \in I}\) be a finite family of elements of \(\hat{\mathbf{A}}(\mathbf{I})\) , let \(v \in \hat{\mathbf{A}}(\mathbf{J})\) and write \(w = v \circ \boldsymbol{u}\) ( \(\S 5\) , no. 1). Then

\[
(v \circ \boldsymbol {u}) ^ {\sim} = \tilde {v} \circ \tilde {\boldsymbol {u}}.\tag{18}
\]

This follows by extending by continuity formula (2) of § 5, no. 1 and from Differentiable and Analytic Manifolds, R, Appendix, no. 6.

In particular we take \(\mathrm{I} = \{\mathrm{U}\}\) , identify A and \(\mathrm{A}^{(\mathrm{U})}\) and consider the images \(\tilde{e}\) and \(\tilde{l}\) of the series \(e(\mathrm{U}) = \sum_{n\geqslant 1}\mathrm{U}^n /n!\) and \(l(\mathrm{U}) = \sum_{n\geqslant 1}(-1)^{n - 1}\mathrm{U}^n /n\) in \(\hat{\mathrm{P}} (\mathrm{A};\mathrm{A})\) . Then \(\| \tilde{\mathbf{U}}^n\| \leqslant 1\) for \(\| x_1\dots x_n\| \leqslant \| x_1\| \dots \| x_n\|\) for \(x_{1},\ldots ,x_{n}\) in A. Therefore the radius of absolute convergence of \(\tilde{e}\) (resp. \(\tilde{l}\) ) is infinite (resp. \(\geqslant 1\) ).

We shall denote by \(e_{A}\) (resp. \(l_{A}\) ) the analytic mapping of A into A (resp. of B into A, where B is the open unit ball of A) defined by the convergent series \(\tilde{e}\) (resp. \(\tilde{l}\) ) and we shall write \(\exp_{\mathbf{A}}(x)=1+e_{\mathbf{A}}(x)\) (for \(x\in A\) ) and

\[
\log_ {\mathrm{A}} (x) = l _ {\mathrm{A}} (x - 1)
\]

(for \(x \in \mathbf{A}\) , \(\| x - 1\| < 1\) ). Then

\[
\exp_ {\mathrm{A}} x = \sum_ {n \geqslant 0} \frac {x ^ {n}}{n !} (x \in \mathrm{A})\tag{19}
\]

\[
\log_ {\mathrm{A}} x = \sum_ {n \geqslant 1} (- 1) ^ {n - 1} \frac {(x - 1) ^ {n}}{n} (x \in \mathrm{A}, \| x - 1 \| <   1).\tag{20}
\]

As \((e\circ l)(\mathrm{U}) = (l\circ e)(\mathrm{U}) = \mathrm{U}\) (cf.~§ 6, no. 1), by (18) \(\tilde{e}\circ \tilde{l} = \tilde{l}\circ \tilde{e} = \mathrm{Id}_{\mathbf{A}}\) Therefore (Differentiable and Analytic Manifolds, R, 3.1.9)

\begin{enumerate}
\def\labelenumi{(\arabic{enumi})}
\setcounter{enumi}{20}
\tightlist
\item
\end{enumerate}

\[
\exp_ {\mathrm{A}} (\log_ {\mathrm{A}} (x)) = x \quad (x \in \mathrm{A}, \| x - 1 \| \leqslant 1)\tag{21}
\]

\[
\log_ {\mathbf {A}} (\exp_ {\mathbf {A}} (x)) = x \quad (x \in \mathrm{A}, \| x \| <   \log 2)\tag{22}
\]

for \(\| x\| < \log 2\) implies \(\| \exp_{\mathbf{A}}(x) - 1\| \leqslant \exp \| x\| -1 < 1.\)

Finally we consider A as a complete normed Lie algebra. Then

\[
\| [ x, y ] \| = \| x y - y x \| \leqslant 2 \| x \|. \| y \|.
\]

Proposition 1 of no. 2 implies that the domain of absolute convergence of the formal power series \(\tilde{\mathbf{H}}\) contains the set

\[
\Omega = \{(x, y) \in \mathrm{A} \times \mathrm{A} | \| x \| + \| y \| <   \frac {1}{2} \log 2 \}.
\]

Hence \(\tilde{\mathbf{H}}\) defines an analytic function \(h\colon \Omega \to \mathrm{A}\) . Then \(h(x,y) = \sum_{r,s}\mathrm{H}_{r,s}(x,y)\) (cf.~§ 3, no. 1, Remark 4).

PROPOSITION 3. For \(\| x\| + \| y\| < \frac{1}{2}\log 2,\)

\[
\exp_ {A} x. \exp_ {A} y = \exp_ {A} h (x, y).\tag{23}
\]

It follows from (18) and the relation \(e^{\mathrm{U}}e^{\mathrm{V}} = e^{\mathrm{H}(\mathrm{U},\mathrm{V})}\) that

\[
m \circ (1 + \tilde {e}, 1 + \tilde {e}) = (1 + \tilde {e}) \circ \tilde {\mathrm{H}}
\]

in \(\hat{\mathrm{P}} (\mathrm{A}\times \mathrm{A};\mathrm{A})\) . We therefore deduce from Differentiable and Analytic Manifolds, R, 3.1.9 that (23) is true for \((x,y)\) sufficiently close to \((0,0)\) , whence the proposition follows by analytic continuation (Differentiable and Analytic Manifolds, R, 3.2.5).

\section{CONVERGENCE OF THE HAUSDORFF SERIES (ULTRAMETRIC CASE)}

In this paragraph we assume that K is a non-discrete complete valued field of characteristic zero, with an ultrametric absolute value. We denote by p the characteristic of the residue field of K (Commutative Algebra, Chapter VI, § 3, no. 2).

If \(p \neq 0\) , we write \(a = |p|\) ; we know (Commutative Algebra, Chapter VI, §6, nos. 2 and 3) that \(0 < a < 1\) and that there exists one and only one valuation \(v\) on \(\mathbf{K}\) with values in \(\mathbf{R}\) whose restriction to \(\mathbf{Q}\) is the \(p\) -adic valuation \(v_p\) and which is such that \(|x| = a^{v(x)}\) for all \(x \in \mathbf{K}\) . Also we write:

\[
\theta = \frac {1}{p - 1}.\tag{1}
\]

If \(p = 0\) , we denote by \(a\) a real number such that \(0 < a < 1\) and by \(v\) a valuation on \(\mathbf{K}\) with values in \(\mathbf{R}\) such that \(|x| = a^{v(x)}\) for all \(x \in \mathbf{K}\) (loc. cit.). Then \(v(x) = 0\) for \(x \in \mathbf{Q}^*\) . Also we write:

\[
\theta = 0.\tag{2}
\]

\subsection{p-ADIC UPPER BOUNDS OF THE SERIES exp, log AND H}

In this no. we assume that \(p \neq 0\) .

Lemma 1. Let n be an integer \(\geqslant0\) and let \(n=n_{0}+n_{1}p+\cdots+n_{k}p^{k}\) , with \(0\leqslant n_{i}\leqslant p-1\) , be the p-adic expansion of n.~Let \(\mathrm{S}(n)=n_{0}+n_{1}+\cdots+n_{k}\) . Then

\[
v _ {p} (n!) = \frac {n - \mathrm{S} (n)}{p - 1}.\tag{3}
\]

\(v_{p}(n!) = \sum_{i=1}^{n} v_{p}(i)\) and the number of integers \(i\) between 1 and \(n\) for which \(v_{p}(i) \geqslant j\) is equal to the integral part \([n/p^{j}]\) of \(n/p^{j}\) . Then

\[
v _ {p} (n!) = \sum_ {j \geqslant 0} j ([ n / p ^ {j} ] - [ n / p ^ {j + 1} ]) = \sum_ {j \geqslant 1} [ n / p ^ {j} ].
\]

As \([n / p^j] = \sum_{i\geqslant j}n_i p^{i - j}\) , the lemma follows.

Lemma 2. \(v(n) \leqslant v(n!) \leqslant (n - 1)\theta\) and \(v(n) \leqslant (\log n) / (\log p)\) for every integer \(n \geqslant 1\) .

\(v(n!) = v_{p}(n!) = (n - S(n))\theta \leqslant (n - 1)\theta\) by Lemma 1.

On the other hand, \(n \geqslant p^{v(n)}\) , whence \(v(n) \leqslant (\log n) / (\log p)\) .

Let \(\mathbf{I} = \{\mathbf{U},\mathbf{V}\}\) be a set of two elements and let

\[
\mathrm{H} = \sum_ {r, s \geqslant 0} \mathrm{H} _ {r, s} (\mathrm{U}, \mathrm{V}) \in \hat {\mathrm{L}} _ {\mathbf {Q}} (\mathrm{I})
\]

be the Hausdorff series (§ 6, no. 4, Definition 1). Let \(\mathbf{Z}_{(p)}\) be the local ring of \(\mathbf{Z}\) relative to the prime ideal \((p)\) and \((e_b)_{b \in \mathbf{B}}\) a basis of \(\mathrm{L}_{\mathbf{Z}_{(p)}}(\mathrm{I})\) over \(\mathbf{Z}\) (§ 2, no. 11, Theorem 1). It is also a basis of \(\mathrm{L}_{\mathbf{Q}}(\mathrm{I})\) over \(\mathbf{Q}\) .

PROPOSITION 1. Let \(r\) and \(s\) be two integers \(\geqslant 0\) . If \(H_{r,s} = \sum_{b \in B} \lambda_b e_b\) , where \(\lambda_b \in \mathbf{Q}\) , is the decomposition of \(H_{r,s}\) with respect to the basis \((e_b)_{b \in B}\) , then

\[
v _ {p} (\lambda_ {b}) \geqslant - (r + s - 1) \theta \quad for all b \in \mathbf {B}.\tag{4}
\]

The ring \(\mathrm{A}_{\mathbf{Z}_{(p)}}(\mathrm{I})\) is identified with the sub- \(\mathbf{Z}_{(p)}\) -module of \(\mathrm{A}_{\mathbf{Q}}(\mathrm{I})\) generated by the words \(w \in \mathrm{Mo}(\mathrm{I})\) . As \(\mathrm{L}_{\mathbf{Z}_{(p)}}(\mathrm{I})\) is a direct factor of \(\mathrm{A}_{\mathbf{Z}_{(p)}}(\mathrm{I})\) ,

\[
\mathrm{L} _ {\mathbf {Z} _ {(p)}} (\mathrm{I}) = \mathrm{A} _ {\mathbf {Z} _ {(p)}} (\mathrm{I}) \cap \mathrm{L} _ {\mathbf {Q}} (\mathrm{I}).\tag{5}
\]

Let \(f\) be the integer such that \(f \leqslant (r + s - 1)\theta < f + 1\) . Relation (4) is equivalent to \(v_{p}(\lambda_{b}) \geqslant -f\) for all \(b \in B\) , that is \(H_{r,s} \in p^{-f}L_{\mathbf{Z}_{(p)}}(I)\) . But this is equivalent also by (5) to \(H_{r,s} \in p^{-f}A_{\mathbf{Z}_{(p)}}(I)\) .

By formula (11) of §6, no. 4, it suffices to show that, for every integer \(m \geqslant 1\) and all integers \(r_1, \ldots, r_m, s_1, \ldots, s_m\) such that

\[
r _ {1} + \dots + r _ {m} = r, \qquad s _ {1} + \dots + s _ {m} = s,\tag{6}
\]

\[
r _ {i} + s _ {i} \geqslant 1 \quad \text { for } 1 \leqslant i \leqslant m,
\]

we have

\[
v _ {p} (m. r! \dots r _ {m}! s _ {1}! \dots s _ {m}!) \leqslant f.\tag{7}
\]

But by Lemma 2, \(v_{p}(r_{i}!s_{i}!) \leqslant (r_{i} + s_{i} - 1)\theta\) and \(v_{p}(m) \leqslant v_{p}(m!) \leqslant (m - 1)\theta\) ; the left hand side of (7) is therefore bounded above by

\[
\theta (m - 1 + \sum_ {i = 1} ^ {m} \left(r _ {i} + s _ {i} - 1\right)) = \theta (r + s - 1);
\]

as it is an integer, it is \(\leqslant f\) , which completes the proof.

\subsection{NORMED LIE ALGEBRAS}

DEFINITION 1. A normed Lie algebra over \(\mathbf{K}\) is a Lie algebra with a norm such that

\begin{enumerate}
\def\labelenumi{(\arabic{enumi})}
\setcounter{enumi}{7}
\tightlist
\item
\end{enumerate}

\[
\| x + y \| \leqslant \sup (\| x \|, \| y \|)\tag{8}
\]

\[
\| [ x, y ] \| \leqslant \| x \|. \| y \|\tag{9}
\]

for all x, y in g.

Throughout the rest of this paragraph, \(\mathfrak{g}\) denotes a complete normed Lie algebra.

For every finite set I we define as in § 7, no. 1 a continuous Lie algebra homomorphism \(u \mapsto \tilde{u}\) of \(\hat{\mathbf{L}}(\mathbf{I})\) into \(\hat{\mathrm{P}}(\mathfrak{g}^{\mathrm{I}}; \mathfrak{g})\) . We see as in § 7 that if \(u = \sum_{\nu} u_{\nu}\) , with \(u_{\nu} \in \mathbf{L}^{\nu}(\mathbf{I})\) for \(\nu \in \mathbf{N}^{\mathbf{I}}\) , then \(\tilde{u} = \sum_{\nu} \tilde{u}_{\nu}\) , where \(\tilde{u}_{\nu}\) is the polynomial mapping \((t_i)_{i \in \mathbf{I}} \mapsto u_{\nu}(t_i)\) defined in § 2, no. 4. The composition formula (2) of § 7, no. 1, remains valid.

\subsection{GROUP DEFINED BY A COMPLETE NORMED LIE ALGEBRA}

Let \(H = \sum_{r,s \geq 0} H_{r,s} \in \hat{L}(\{U, V\})\) be the Hausdorff series (§ 6, no. 4, Definition 1). We shall show that the corresponding formal power series with continuous components

\[
\tilde {\mathrm{H}} = \sum_ {r, s \geqslant 0} \tilde {\mathrm{H}} _ {r, s} \in \hat {\mathrm{P}} (\mathfrak {g} \times \mathfrak {g}, \mathfrak {g})\tag{10}
\]

is convergent (Differentiable and Analytic Manifolds, R, 4.1.1).

Let \(r \geqslant 0\) , \(s \geqslant 0\) be such that \(r + s \neq 0\) and let \(\| \tilde{\mathrm{H}}_{r,s} \|\) be the norm of the continuous polynomial \(\tilde{\mathrm{H}}_{r,s}\) (Differentiable and Analytic Manifolds, R, Appendix, no. 2).

Lemma 3.

\[
\left\| \tilde {\mathrm{H}} _ {r, s} \right\| \leqslant a ^ {- (r + s - 1) \theta}.
\]

Let B be a Hall set relative to I and let \(H_{r,s} = \sum_{b \in B} \lambda_b e_b\) be the decomposition of \(H_{r,s}\) with respect to the corresponding basis of \(L(\{U, V\})\) . Then

\[
\left| \lambda_ {b} \right| \leqslant a ^ {- (r + s - 1) \theta}.\tag{11}
\]

This is trivial for \(p = 0\) , for \(\lambda_b \in \mathbf{Q}\) ; and it follows from Proposition 1 of no. 1 for \(p \neq 0\) .

Moreover,

\[
\left\| \tilde {e} _ {b} \right\| \leqslant 1 \quad \text { for } b \in \mathbf {B}.\tag{12}
\]

We show more generally by induction on \(n\) that, for every alternant \(b\) of degree \(n\) in the two indeterminates U and V (§2, no. 6), \(\| \tilde{b} \| \leqslant 1\) . If \(n = 1\) , \(\tilde{b}\) is one of the projections of \(\mathfrak{g} \times \mathfrak{g}\) onto \(\mathfrak{g}\) and hence is of norm \(\leqslant 1\) ; if \(n > 1\) , there exist two alternants \(b_1\) and \(b_2\) of degrees \(< n\) such that \(b = [b_1, b_2]\) . As the mapping \(\gamma: (x, y) \mapsto [x, y]\) of \(\mathfrak{g} \times \mathfrak{g}\) into \(\mathfrak{g}\) is bilinear of norm \(\leqslant 1\) , we have (Differentiable and Analytic Manifolds, R, Appendix, no. 4)

\[
\| \tilde {b} \| = \| \gamma \circ (\tilde {b} _ {1}, \tilde {b} _ {2}) \| \leqslant \| \tilde {b} _ {1} \|. \| \tilde {b} _ {2} \| \leqslant 1.\tag{13}
\]

Relations (11) and (12) imply the lemma.

PROPOSITION 2. The formal power series \(\tilde{\mathbf{H}}\) is a convergent series (Differentiable and Analytic Manifolds, R, 4.1.1). If \(\mathbf{G}\) is the ball \(\{x\in\mathfrak{g}\mid\|x\|<a^{\theta}\}\) , the domain of absolute convergence of \(\tilde{\mathbf{H}}\) (Differentiable and Analytic Manifolds, R, 4.1.3) contains \(\mathbf{G}\times\mathbf{G}\) .

If \(u\) and \(v\) are two real numbers \(>0\) such that \(u < a^{\theta}\) and \(v < a^{\theta}\) , then (Lemma 3)

\[
\| \tilde {\mathrm{H}} _ {r, s} \| u ^ {r} v ^ {s} \leqslant a ^ {\theta} (u a ^ {- \theta}) ^ {r} (v a ^ {- \theta}) ^ {s}\tag{14}
\]

and \(\| \tilde{\mathrm{H}}_{r,s}\| u^r v^s\) tends to 0 when \(r + s\) tends to infinity.

We denote by \(h \colon G \times G \to \mathfrak{g}\) the analytic function (Differentiable and Analytic Manifolds, R, 4.2.4) defined by \(\tilde{\mathbf{H}}\) , that is by the formula

\[
h (x, y) = \sum_ {r, s \geqslant 0} \tilde {\mathrm{H}} _ {r, s} (x, y) = \sum_ {r, s \geqslant 0} \mathrm{H} _ {r, s} (x, y) \quad \text { for } (x, y) \in \mathrm{G} \times \mathrm{G}.\tag{15}
\]

This function is called the Hausdorff function of g.

Let \((x,y)\in \mathbf{G}\times \mathbf{G}\) . Then

\begin{enumerate}
\def\labelenumi{(\arabic{enumi})}
\setcounter{enumi}{15}
\tightlist
\item
\end{enumerate}

\[
\left\| \tilde {\mathrm{H}} _ {r, s} (x, y) \right\| \leqslant \sup (\left\| x \right\|, \left\| y \right\|)\tag{16}
\]

\[
\| h (x, y) \| \leqslant \sup (\| x \|, \| y \|).\tag{17}
\]

\begin{enumerate}
\def\labelenumi{(\arabic{enumi})}
\setcounter{enumi}{16}
\tightlist
\item
  follows immediately from (16) and (16) is trivial for \(r = s = 0\) ; if \(r \geqslant 1\) , then
\end{enumerate}

\[
\begin{array}{r l} \| \tilde {\mathrm{H}} _ {r, s} (x, y) \| & \leqslant \| \tilde {\mathrm{H}} _ {r, s} \| \| x \| ^ {r} \| y \| ^ {s} \\ & \leqslant \| x \| \left(\frac {\| x \|}{a ^ {\theta}}\right) ^ {r - 1} \left(\frac {\| y \|}{a ^ {\theta}}\right) ^ {s} \\ & \leqslant \| x \|; \end{array}
\]

we argue similarly if \(s \geqslant 1\) .

In particular, \(\| h(x,y)\| < a^{\theta}\) for \((x,y)\in \mathbf{G}\times \mathbf{G}\) .

PROPOSITION 3. Let G be the ball \(\{x\in\mathfrak{g}\mid\|x\|<a^{\theta}\}\) . The analytic mapping

\[
h \colon \mathrm{G} \times \mathrm{G} \rightarrow \mathrm{G}
\]

makes G into a group in which 0 is the identity element and -x is the inverse of x for all \(x \in G\) . Moreover, if R is a real number such that \(0 < R < a^{\theta}\) , the ball

\[
\{x \in \mathfrak {g} | \| x \| <   \mathrm{R} \}
\]

(resp. \(\{x \in \mathfrak{g} \mid \|x\| \leqslant R\}\)) is an open subgroup of G.

As \(\mathrm{H}(\mathrm{U}, -\mathrm{U}) = 0\) and \(\mathrm{H}(0, \mathrm{U}) = \mathrm{H}(\mathrm{U}, 0) = \mathrm{U}\) , \(h(x, -x) = 0\) and

\[
h (0, x) = h (x, 0) = x
\]

for all \(x \in G\) . It therefore remains to prove the associativity formula

\[
h (h (x, y), z) = h (x, h (y, z)) \quad \text { for } x, y, z \text { in } G.\tag{18}
\]

As

\[
\mathrm{H} (\mathrm{H} (\mathrm{U}, \mathrm{V}), \mathrm{W}) = \mathrm{H} (\mathrm{U}, \mathrm{H} (\mathrm{V}, \mathrm{W}))
\]

in \(\hat{\mathbf{L}} (\{\mathrm{U},\mathrm{V},\mathrm{W}\})\) (§ 6, no. 5, Proposition 4), we have

\[
\tilde {\mathrm{H}} \circ (\tilde {\mathrm{H}} \times \mathrm{Id} _ {\mathfrak {g}}) = \tilde {\mathrm{H}} \circ (\mathrm{Id} _ {\mathfrak {g}} \times \tilde {\mathrm{H}})\tag{19}
\]

in \(\hat{\mathbf{P}} (\mathfrak{g}\times \mathfrak{g}\times \mathfrak{g};\mathfrak{g})\) (no. 2) and (19) implies (18) by (16) and Differentiable and Analytic Manifolds, R, 4.1.5.

*In other words (Chapter III, §1), G with the Hausdorff function is a Lie group.*

\subsection{EXPONENTIAL IN COMPLETE NORMED ASSOCIATIVE ALGEBRAS}

In this no. A will denote a unital associative algebra with a norm \(x \mapsto \| x \|\) satisfying the conditions:

\[
\begin{array}{c} \| x + y \| \leqslant \sup (\| x \|, \| y \|) \\ \| x y \| \leqslant \| x \|. \| y \| \\ \| 1 \| = 1 \end{array}
\]

for \(x, y\) in A, and complete with this norm. The results of the second and third paragraphs of § 7, no. 3, remain valid.

We take \(\mathbf{I} = \{\mathrm{U}\}\) and consider the images \(\tilde e\) and \(\tilde{l}\) of the series \(e(\mathrm{U}) = \sum_{n\geqslant 1}\frac{\mathrm{U}^n}{n!}\) and \(l(\mathrm{U}) = \sum_{n\geqslant 1}(-1)^{n - 1}\frac{U^n}{n}\) in \(\hat{\mathrm{P}} (\mathrm{A};\mathrm{A})\) . Then:

\begin{enumerate}
\def\labelenumi{(\arabic{enumi})}
\setcounter{enumi}{19}
\tightlist
\item
\end{enumerate}

\[
\left\| \left(\frac {\mathrm{U} ^ {n}}{n !}\right) ^ {\sim} \right\| \leqslant a ^ {- (n - 1) \theta}\tag{20}
\]

\[
\left\| \left(\frac {\mathrm{U} ^ {n}}{n}\right) ^ {\sim} \right\| \leqslant a ^ {- \frac {\log n}{\log p}}\tag{21}
\]

by Lemma 2 of no. 1. Hence the radius of absolute convergence of the series \(\tilde{e}\) (resp. \(\tilde{l}\) ) is \(\geqslant a^\theta\) (resp. \(\geqslant 1\) ) (Differentiable and Analytic Manifolds, R, 4.1.3). For \(R > 0\) , let \(G_R = \{x \in A \mid \|x\|<R\}\) ; we write \(G = G_\theta\) . The series \(\tilde{e}\) (resp. \(\tilde{l}\) ) defines an analytic mapping \(e_{A}\) (resp. \(l_{A}\) ) of \(G\) (resp. \(G_{1}\) ) into A. We write:

\begin{enumerate}
\def\labelenumi{(\arabic{enumi})}
\setcounter{enumi}{21}
\tightlist
\item
\end{enumerate}

\[
\exp_ {\mathrm{A}} (x) = 1 + e _ {\mathrm{A}} (x) = \sum_ {n \geqslant 0} \frac {x ^ {n}}{n !}\tag{22}
\]

\[
\log_ {\mathrm{A}} (x) = l _ {\mathrm{A}} (x - 1) = \sum_ {n \geqslant 1} (- 1) ^ {n - 1} \frac {(x - 1) ^ {n}}{n} \quad \text { for } x - 1 \in \dot {\mathrm{G}} _ {1}\tag{23}
\]

(we omit the index A when no confusion can arise). For \(x \in G_R\) and \(n \geqslant 1\) ,

\[
\left\| \frac {x ^ {n}}{n} \right\| \leqslant \left\| \frac {x ^ {n}}{n !} \right\| <   R ^ {n} a ^ {- (n - 1) \theta} = R \left(\frac {R}{a ^ {\theta}}\right) ^ {n - 1}\tag{24}
\]

and hence \(e_{\mathrm{A}}(\mathrm{G}_{\mathrm{R}}) \subset \mathrm{G}_{\mathrm{R}}, l_{\mathrm{A}}(\mathrm{G}_{\mathrm{R}}) \subset \mathrm{G}_{\mathrm{R}}\) for \(R \leqslant a^{\theta}\) .

PROPOSITION 4. Let R be a real number such that \(0 < R \leqslant a^{\theta}\) . The mapping \(\exp_{\mathbf{A}}\) defines an analytic isomorphism of \(\mathbf{G}_{R}\) onto \(1 + \mathbf{G}_{R}\) and the inverse isomorphism is the restriction of \(\log_{\mathbf{A}}\) to \(1 + \mathbf{G}_{R}\) .

\(e(l(\mathbf{X})) = l(e(\mathbf{X})) = \mathbf{X}\) . By (20), (21) and Differentiable and Analytic Manifolds, R, 4.1.5, we deduce that \(e_{\mathrm{A}}(l_{\mathrm{A}}(x)) = l_{\mathrm{A}}(e_{\mathrm{A}}(x))\) for \(x \in \mathrm{G}_{\mathrm{R}}\) . Then

\[
\begin{array}{l l} \exp_ {\mathrm{A}} (\log_ {\mathrm{A}} x) = x & \text {for} x \in 1 + \mathrm{G} _ {\mathrm{R}} \\ \log_ {\mathrm{A}} (\exp_ {\mathrm{A}} x) = x & \text {for} x \in \mathrm{G} _ {\mathrm{R}} \end{array}
\]

which completes the proof.

If A is given the bracket \([x,y]=xy-yx\) , A becomes a complete normed Lie algebra, for \(\|xy-yx\|\leqslant\sup(\|xy\|,\|yx\|)\leqslant\|x\|. \|y\|\) . Proposition 2 of no. 3 implies that the domain of absolute convergence of \(\tilde{H}\) contains \(G\times G\) and \(\tilde{H}\) therefore defines an analytic function \(h:G\times G\to A\) ; then

\[
h (x, y) = \sum_ {r, s \geqslant 0} \mathrm{H} _ {r, s} (x, y).\tag{25}
\]

PROPOSITION 5. For x, y in G,

\[
\exp_ {\mathrm{A}} x. \exp_ {\mathrm{A}} y = \exp_ {\mathrm{A}} h (x, y).\tag{26}
\]

\(e^{\mathrm{U}}e^{\mathrm{V}} = e^{\mathrm{H(U,V)}}\) and hence

\[
m \circ (1 + \tilde {e}, 1 + \tilde {e}) = (1 + \tilde {e}) \circ \tilde {\mathrm{H}}
\]

in \(\hat{\mathrm{P}}(\mathbf{A} \times \mathbf{A}; \mathbf{A})\) (where m denotes multiplication on A). The proposition then follows from Proposition 2, Lemma 3 and Differentiable and Analytic Manifolds, R, 4.1.5.


\appendixsection{Möbius function}

Let \(n\) be an integer \(\geqslant 1\) . If \(n\) is divisible by the square of a prime number, we write \(\mu(n) = 0\) . If \(n\) is not divisible by the square of a prime number, we write \(\mu(n) = (-1)^k\) , where \(k\) is the number of prime divisors of \(n\) . The function \(\mu: \mathbf{N}^* \to \{-1, 0, 1\}\) thus defined is called the Möbius function.

Recall that given two integers \(n_1 \geqslant 1, n_2 \geqslant 1\) , we write \(n_1|n_2\) if \(n_1\) divides \(n_2\) .

PROPOSITION. (i) The function \(\mu\) is the unique mapping of \(\mathbf{N}^*\) into \(\mathbf{Z}\) such that \(\mu(1) = 1\) and

\[
\sum_ {d \mid n} \mu (d) = 0\tag{1}
\]

for every integer \(n > 1\) .

\begin{enumerate}
\def\labelenumi{(\roman{enumi})}
\setcounter{enumi}{1}
\tightlist
\item
  Let \(s\) and \(t\) be two mappings of \(\mathbf{N}^*\) into a commutative group written additively. In order that
\end{enumerate}

\[
s (n) = \sum_ {d | n} t (d) \quad \text { for   every   integer } n \geqslant 1,\tag{2}
\]

it is necessary and sufficient that

\[
t (n) = \sum_ {d \mid n} \mu (d) s \left(\frac {n}{d}\right) \quad for every integer n \geqslant 1.\tag{3}
\]

The uniqueness assertion in (i) is obvious, for (1) allows us to determine \(\mu(n)\) by induction on \(n\) . We show that the function \(\mu\) satisfies (1). Let \(n\) be an integer \(>1\) . Let P be the set of prime divisors of \(n\) and let \(n = \prod_{p \in \mathbb{P}} p^{\nu_p(n)}\) be the decomposition of \(n\) into prime factors. If \(d\) is a divisor of \(n\) , then \(\mu(d) = 0\) unless \(d\) is of the form \(\prod_{p \in \mathbb{H}} p\) , where H is a subset of P. Then

\[
\begin{array}{r l} \sum_ {d | n} \mu (d) & = \sum_ {\mathrm{H} \subset \mathrm{P}} (- 1) ^ {\mathrm{CardH}} \\ & = \sum_ {k = 0} ^ {\mathrm{CardP}} \left(\frac {\mathrm{CardP}}{k}\right) (- 1) ^ {k} = (1 - 1) ^ {\mathrm{CardP}} = 0. \end{array}
\]


Let \(s\) and \(t\) be two mapping of \(\mathbf{N}^*\) into a commutative group written additively. Let \(n \in \mathbf{N}^*\) . If (2) holds, then

\[
\begin{array}{r l} \sum_ {d | n} \mu (d) s \left(\frac {n}{d}\right) & = \sum_ {d | n} \mu (d) \sum_ {\delta \mid (n / d)} t (\delta) = \sum_ {d \delta | n} \mu (d) t (\delta) \\ & = \sum_ {\delta | n} t (\delta) \sum_ {d \mid (n / \delta)} \mu (d) = t (n). \end{array}
\]

Conversely, if (3) holds, then

\[
\sum_ {d \mid n} t (d) = \sum_ {d \mid n} \sum_ {\delta \mid d} \mu (\delta) s \left(\frac {d}{\delta}\right) = \sum_ {d \mid n} s (d) \sum_ {\delta \mid (n / d)} \mu (\delta) = s (n),
\]

which completes the proof.

Formula (3) is called the Möbius inversion formula.


\specialchapter{Exercises}

\exercisesfor{1}

\begin{enumerate}
\def\labelenumi{\arabic{enumi}.}
\item
  Suppose that K is a field of characteristic 0. Let E be a cocommutative bigebra of finite rank over K. Show that \(P(E) = \{0\}\) (apply Theorem 1 of no. 6).
\item
  Let E be a bigebra such that \(\mathrm{P}(\mathrm{E})=\{0\}\) and let \((\mathrm{E}_{n})_{n\geqslant0}\) be a filtration of E compatible with its bigebra structure. Show by induction on n that \(E_{n}^{+}=\{0\}\) for all \(n\geqslant0\) and deduce that \(E^{+}=\{0\}\) , i.e.~that E reduces to K.
\item
  Let \(G\) be a monoid and \(E = K[G]\) its algebra (Algebra, Chapter III, §2, no. 6).
\end{enumerate}

\begin{enumerate}
\def\labelenumi{(\alph{enumi})}
\item
  Show that there exists on E one and only one cogebra structure such that \(c(g) = g \otimes g\) for all \(g \in G\) ; this structure is compatible with the algebra structure on E and makes E into a cocommutative bigebra whose counit \(\varepsilon\) is such that \(\varepsilon(g) = 1\) for all \(g \in G\) .
\item
  Show that every primitive element of E is zero. Deduce that, if \(G \neq \{e\}\) , the bigebra E admits no filtration compatible with its bigebra structure (apply Exercise 2).
\item
  Suppose that \(\mathbf{K}\) is an integral domain. Show that the elements of \(\mathrm{G}\) are the only non-zero elements \(x\in \mathrm{E}\) such that \(c(x) = x\otimes x\) . Show that \(\mathrm{E}\) has an inversion \(i\) (cf.~Algebra, Chapter III, § 11, Exercise 4) if and only if \(\mathrm{G}\) is a group and that in that case \(i(g) = g^{-1}\) for all \(g\in \mathrm{G}\) .
\end{enumerate}

\begin{enumerate}
\def\labelenumi{\arabic{enumi}.}
\setcounter{enumi}{3}
\item
  Let \(\mathbf{E}\) be a cocommutative bigebra and let \(\mathbf{G}\) be the set of \(g \in \mathbf{E}\) such that \(\varepsilon(g) = 1\) and \(c(g) = g \otimes g\) . Show that \(\mathbf{G}\) is stable under multiplication and that it is a group if \(\mathbf{E}\) has an inversion. Show that, if \(\mathbf{K}\) is a field, the elements of \(\mathbf{G}\) are linearly independent over \(\mathbf{K}\) .
\item
  Let E be a bigebra and let \(E' = \text{Hom}(E, K)\) be its dual with the algebra structure derived by duality from the cogebra structure on E (Algebra, Chapter III, §11, no. 1). Let m be the kernel of the homomorphism \(u \mapsto u(1)\) of \(\mathbf{E}'\) onto \(\mathbf{K}\) ; it is an ideal of \(\mathbf{E}'\) . Show that, if \(u \in \mathfrak{m}^2\) and \(x \in \mathrm{P}(\mathrm{E})\) , then \(u(x) = 0\) . When \(\mathbf{K}\) is a field show that \(\mathrm{P}(\mathrm{E})\) is the orthogonal of \(\mathfrak{m}^2\) in \(\mathbf{E}^+\) ; deduce that \(\dim \mathrm{P}(\mathrm{E}) \leqslant \dim \mathfrak{m} / \mathfrak{m}^2\) and that we have equality if \(\mathbf{E}\) is of finite rank over \(\mathbf{K}\) .
\item
  Let E be a bigebra and let \(u \in \mathrm{E}' = \operatorname{Hom}(\mathrm{E}, \mathrm{K})\) (cf.~Exercise 5). Let \(g \in \mathrm{E}\) be such that \(\varepsilon(g) = 1\) and \(c(g) = g \otimes g\) . Show that the mapping \(u \mapsto u(g)\) is an algebra homomorphism of \(\mathrm{E}'\) into \(\mathrm{K}\) and that the mapping \(y \mapsto gy\) is an endomorphism of the cogebra \(\mathrm{E}\) ; show that this endomorphism is an automorphism if \(\mathrm{E}\) has an inversion.
\item
  Suppose that K is a field. Let E be a bigebra satisfying the following conditions:
\end{enumerate}

\begin{enumerate}
\def\labelenumi{(\roman{enumi})}
\item
  E is cocommutative;
\item
  E has an inversion (Algebra, Chapter III, § 11, Exercise 4);
\item
  \(\mathrm{P(E)} = \{0\}\) ;
\item
  E is of finite rank over K.
\end{enumerate}

Let \(\mathbf{E}'\) denote the dual K-algebra of the cogebra E. Let G denote the set of elements \(g \in \mathbf{E}\) such that \(\varepsilon(g) = 1\) and \(c(g) = g \otimes g\) ; it is a group (cf.~Exercise 4).

\begin{enumerate}
\def\labelenumi{(\alph{enumi})}
\item
  Let \(g \in \mathbf{G}\) and let \(\mathfrak{m}_g\) be the kernel of the homomorphism \(u \mapsto u(g)\) of \(\mathbf{E}'\) into \(\mathbf{K}\) . Show that \(\mathfrak{m}_g = \mathfrak{m}_g^2\) (reduce it to the case \(g = 1\) using Exercise 6; then apply Exercise 5).
\item
  Suppose that K is algebraically closed. Show that the ideals \(m_{g}\) are the only maximal ideals of \(E'\) ; deduce that their intersection is \(\{0\}\) , that \(E'\) is identified with the product \(K^{G}\) and that E is identified with the bigebra K[G] of the finite group G (cf.~Exercise 3).
\end{enumerate}

\begin{enumerate}
\def\labelenumi{\arabic{enumi}.}
\setcounter{enumi}{7}
\item
  Show that the enveloping bigebra of a Lie algebra \(\mathfrak{g}\) has an inversion \(i\) and that \(i(x) = -x\) for all \(x \in \sigma(\mathfrak{g})\) .
\item
  Suppose that K is a Q-algebra. Let g be a Lie algebra admitting a K-basis and let U be its enveloping bigebra with the canonical filtration \((\mathrm{U}_{n})_{n\geqslant0}\) . Let \(x\in U^{+}\) and let n be an integer \(\geqslant1\) . Show that x belongs to \(U_{n}^{+}\) if and only if \(c^{+}(x)\) belongs to \(\sum_{\substack{i+j=n\\ i,j\geqslant1}}\mathrm{Im}(\mathrm{U}_{i}^{+}\otimes\mathrm{U}_{j}^{+})\) .
\item
  Suppose that K is a Q-algebra. Let E be a cocommutative K-bigebra with a filtration compatible with its bigebra structure. Show that the morphism
\end{enumerate}

\[
f _ {\mathbf {E}} \colon \mathrm{U} (\mathrm{P} (\mathrm{E})) \rightarrow \mathrm{E}
\]

defined in no. 4, Proposition 8 is an isomorphism, if P(E) is assumed to be a free K-module (same proof as for Theorem 1).

¶ 11. Let E be a bigebra and let I be a finite set. We are given a basis \((e_{\alpha})\)

of E indexed by the elements \(\alpha\) of \(\mathbf{N}^{\mathrm{I}}\) and such that:

\begin{enumerate}
\def\labelenumi{(\roman{enumi})}
\item
  \(e_{0} = 1\) ;
\item
  \(c(e_{\gamma}) = \sum_{\alpha + \beta = \gamma} e_{\alpha} \otimes e_{\beta}\) for all \(\gamma \in \mathbf{N}^{\mathrm{I}}\) .
\end{enumerate}

Condition (ii) implies that E is isomorphic as a cogebra to the cogebra \(\mathrm{TS}(\mathrm{K}^{\mathrm{I}})\) of symmetric tensors of \(\mathrm{K}^{\mathrm{I}}\) (cf.~Algebra, Chapter IV, § 5, no. 7).

\begin{enumerate}
\def\labelenumi{(\alph{enumi})}
\tightlist
\item
  Let \(\mathrm{E}' = \mathrm{Hom}(\mathrm{E},\mathrm{K})\) be the dual algebra of \(\mathbf{E}\) and let \(\Lambda = \mathrm{K}[[(\mathrm{X}_i)_{\mathrm{i}\in I}]]\) be the algebra of formal power series in indeterminates \(\mathbf{X}_i\) indexed by I. For all \(u\in \mathrm{E}'\) let \(\phi_u\) be the formal power series
\end{enumerate}

\[
\sum_ {\alpha} u (e _ {\alpha}) x ^ {\alpha},
\]

where \(x^{\alpha} = \prod_{i\in I}\mathrm{X}_{i}^{\alpha (i)}\)

Show that \(u \mapsto \phi_u\) is an isomorphism of \(E'\) onto \(\Lambda\) and that this isomorphism maps the topology of simple convergence on \(E'\) to the product topology on \(\Lambda\) .

\begin{enumerate}
\def\labelenumi{(\alph{enumi})}
\setcounter{enumi}{1}
\tightlist
\item
  Let \(c_{\alpha \beta \gamma}\) be the constants of structure of the algebra \(E\) relative to the basis \((e_{\alpha})\) . Then
\end{enumerate}

\[
e _ {\alpha} e _ {\beta} = \sum_ {\gamma} c _ {\alpha \beta \gamma} e _ {\gamma}.
\]

We write x instead of \((\mathbf{X}_{i})_{i\in\mathbb{I}}\) and y instead of \((\mathbf{Y}_{i})_{i\in\mathbb{I}}\) , where the \(Y_{i}\) are new indeterminates. Show that there exists a family \(f(\boldsymbol{x},\boldsymbol{y})=(f_{i}(\boldsymbol{x},\boldsymbol{y}))_{i\in\mathbb{I}}\) of formal power series in the variables x, y such that

\[
\boldsymbol {f} (\boldsymbol {x}, \boldsymbol {y}) ^ {\gamma} = \sum_ {\alpha , \beta} c _ {\alpha \beta \gamma} \boldsymbol {x} ^ {\alpha} \boldsymbol {y} ^ {\beta} \quad \text { for   all } \gamma \in \mathbf {N} ^ {\mathrm{I}},
\]

where we write, as above, \(x^{\alpha} = \prod_{i} X_{i}^{\alpha(i)}\) and similarly for \(y^{\beta}\) and \(f(x, y)^{\gamma}\) .

Show that \(f(x, y)\) is a formal group law over \(K\) of dimension \(n = \text{Card}(I)\) , in the sense of Chapter I, § 1, Exercise 24†; define a Lie algebra isomorphism of this formal group law onto the Lie algebra \(P(E)\) , which has basis the \(e_{\alpha}\) with \(|\alpha| = 1\) .

\begin{enumerate}
\def\labelenumi{(\alph{enumi})}
\setcounter{enumi}{2}
\item
  Conversely, show that every formal group law over K in the x and y can be obtained by the above procedure, which is unique up to isomorphism.
\item
  When K is a Q-algebra, apply the above to the enveloping bigebra of a Lie algebra g with a finite basis over K. Deduce the existence (and uniqueness, up to isomorphism) of a formal group law with g as Lie algebra.
\item
  Give explicitly the bigebra corresponding to the formal group law in one parameter \(f(x, y) = x + y\) (resp. \(f(x, y) = x + y + xy\) ).
\end{enumerate}

¶ 12. Suppose that K is a field of characteristic p > 0.

\begin{enumerate}
\def\labelenumi{(\alph{enumi})}
\tightlist
\item
  Let \(\mathfrak{g}\) be a Lie \(p\) -algebra (Chapter I, § 1, Exercise 20) and let \(\tilde{\mathbf{U}}\) be its
\end{enumerate}


restricted enveloping bigebra (Chapter I, §2, Exercise 6). Show that there exists on \(\tilde{\mathbf{U}}\) one and only one bigebra structure which is compatible with its algebra structure and for which the elements of \(g\) are primitive. Show that \(\tilde{\mathbf{U}}\) is cocommutative and that \(\mathrm{P}(\tilde{\mathbf{U}}) = g\) .

\begin{enumerate}
\def\labelenumi{(\alph{enumi})}
\setcounter{enumi}{1}
\item
  If \(n\) is an integer \(\geqslant 0\) , let \(\tilde{\mathbf{U}}_n\) denote the vector subspace of \(\tilde{\mathbf{U}}\) generated by the products \(x_1\cdots x_n\) where \(x_i \in \mathfrak{g}\) for all \(i\) . Show that \((\tilde{\mathbf{U}}_n)_{n \geqslant 0}\) is a filtration of \(\tilde{\mathbf{U}}\) compatible with its bigebra structure.
\item
  Let \(\mathbf{E}\) be a bigebra and let \(\mathfrak{g} = \mathrm{P}(\mathbf{E})\) be the Lie algebra of its primitive elements. The mapping \(x\mapsto x^p\) leaves \(\mathfrak{g}\) stable and gives \(\mathfrak{g}\) a Lie \(p\) -algebra structure. Show that the canonical injection \(\mathfrak{g}\to \mathbf{E}\) can be extended uniquely to a bigebra morphism \(\tilde{\mathbf{U}}\rightarrow \mathbf{E}\) and that this morphism is injective (use the fact that, if \((x_{i})_{i\in I}\) is a basis of \(\mathfrak{g}\) with a total ordering, the monomials \(\prod_{i\in I}x_i^{n_i}(0\leqslant n_i < p)\) form a basis of \(\mathbf{U}\) (loc. cit.) and argue as in the proof of Lemma 2).
\end{enumerate}
